\begin{textsample}{NEG Dim 1 – 2018\_c6 – Score -3.00 – 2018\_c6/enade\_2018\_tec\_des\_inter\_0a7fad3a23.txt}  \label{ex:f1_neg_199}
CAPÍTULO 6 ANÁLISE TÉCNICA DA PROVA Este capítulo tem por objetivo apresentar o desempenho dos estudantes concluintes de Tecnologia em Design de Interiores no Enade/2018.
Para isso, foram calculadas as estatísticas básicas da prova em seu todo ( seção 6.1.1 ), bem como as estatísticas relacionadas aos Componentes de Formação Geral ( seção 6.1.2 ) e de Conhecimento Específico da Área ( seção 6.1.3 ).
Dadas as suas características, foram analisadas, em separado, as questões objetivas ( seção 6.2 ) e as questões discursivas ( seção 6.3 ).
Para as questões objetivas, foram disponibilizados os índices de facilidade e de discriminação Ponto-Bisserial, também em separado, para os Componentes de Formação Geral ( seção 6.2.1 ) e de Conhecimento Específico ( seção 6.2.2 ).
De \textbf{cada} componente, uma das questões foi escolhida para exemplificar a análise gráfica, relacionando as alternativas escolhidas pelos estudantes ( inclusive o gabarito ) com o número de acertos no componente.
No Anexo I, apresenta -se a íntegra da análise gráfica para todas as questões objetivas.
Para \textbf{cada} uma das questões discursivas, os conteúdos dos tipos mais comuns de respostas dos estudantes são apresentados e comparados com o padrão de respostas esperado ( ver Anexo VIII com o padrão de respostas ).
Tomando -se como base as duas questões discursivas do Componente de Formação Geral, na seção 6.3.1.6, são apresentados comentários sobre a correção das respostas em relação à Língua Portuguesa.
Nas tabelas, constam as seguintes estatísticas das notas20 : média do desempenho na prova, erro padrão da média, desvio padrão, nota mínima, mediana e nota máxima para \textbf{cada} um de seus componentes.
Tais estatísticas contemplam o total de estudantes concluintes da Área de Tecnologia em Design de Interiores inscritos e presentes à prova do Enade/2018, tendo em vista agregações, ou por Grandes Regiões ou o país como um todo, ou por Categoria Administrativa, Organização Acadêmica da IES e Modalidade de Ensino.
Em relação aos gráficos de distribuição de notas, o intervalo considerado foi de 10 unidades, aberto à esquerda e fechado à direita, com exceção do primeiro intervalo, [ 0 ; 10 ], fechado em ambos os extremos.
Para os gráficos de distribuição das notas das questões discursivas, foram consideradas mais duas categorias : questão em branco e nota zero. 20 Uma definição dessas estatísticas pode ser encontrada no Glossário. 6.1 ESTATÍSTICAS BÁSICAS DA PROVA Esta seção apresenta estatísticas selecionadas e histogramas da nota geral ( 6.1.1 ) e de \textbf{cada} componente : Formação Geral ( 6.1.2 ) e Conhecimento Específico ( 6.1.3 ).
São, também, apresentadas estatísticas selecionadas de subpopulações, caracterizadas por Grande Região, Categoria Administrativa, Organização Acadêmica e Modalidade de Ensino. 6.1.1 Estatísticas Básicas Gerais Na Tabela 6.1, são apresentadas as Estatísticas Básicas da prova, por Grande Região, dos estudantes concluintes de Tecnologia em Design de Interiores.
A Média das notas da prova, como um todo ( nas seções seguintes serão analisados os Componentes de Formação Geral e de Conhecimento Específico ), foi 47,3, sendo que os alunos da região Norte obtiveram a Média mais baixa ( 45,6 ), e os da região Sul obtiveram a Média mais alta ( 48,0 ).
As demais médias foram : 47,9, na região Nordeste, 46,9, na região Sudeste e 46,0, na região Centro-Oeste.
O Desvio padrão para o Brasil, como um todo, foi 12,5, sendo o maior Desvio padrão encontrado na região Norte ( 14,1 ), e o menor, na região Sul ( 12,0 ), indicando uma dispersão um pouco menor das notas dessa última região.
A região que obteve a maior nota Máxima foi a Nordeste ( 81,4 ), ao passo que a região que atingiu a menor nota Máxima foi a Norte ( 70,8 ).
A Mediana do Brasil, como um todo, foi 47,5, sendo a maior Mediana obtida na região Nordeste ( 48,2 ), e a menor obtida na Norte ( 45,8 ).
A nota Mínima do Brasil foi 2,1, e foi obtida na região Sul.
A Mínima na região Norte foi 9,6, a maior nota Mínima dentre as cinco Grandes Regiões, e, nas demais, foram : 7,7 na região Nordeste, 6,4 na Sudeste e 4,6 na Centro-Oeste.
Considerando -se as notas segundo Grande Região, observa -se que não existe diferença estatisticamente significativa ao nível de 95 \% entre as Médias, obtidas nas cinco Grandes Regiões.22. 22 Todas as comparações deste capítulo utilizam os intervalos de 95 \%.
Os erros-padrão da média que possibilitam os testes \textbf{estão} disponíveis nas tabelas.
O comportamento das notas dos estudantes de todo o Brasil pode ser observado no Gráfico 6.1, que apresenta um histograma com a distribuição das mesmas : uma distribuição unimodal com a moda no intervalo ( 40 ; 50 ], seguido de perto pelo intervalo ( 50 ; 60 ].
Na Tabela 6.2, são apresentadas informações referentes à Média da nota final, desagregadas por Categoria Administrativa, por Organização Acadêmica e por Modalidade de Ensino.
A Média das notas da prova, como um todo, foi 47,4.
Em relação à Categoria Administrativa, os estudantes das IES Públicas obtiveram Média mais alta ( 53,1 ) que a Média nacional, e os das IES Privadas obtiveram Média mais baixa ( 46,9 ).
Observa -se que existe diferença estatisticamente significativa dentre as médias das notas das IES Públicas e as das Privadas.
A diferença entre as médias das IES Públicas e as das Privadas ( 6,2 ) é superior à diferença entre as médias das regiões Sul e Norte ( 2,4 ), a maior e a menor Média, caracterizando -se uma maior diversidade administrativa do que regional.
O Desvio padrão para as IES Públicas ( 11,8 ) e para IES Privadas ( 12,4 ) foi inferior ao do Brasil, como um todo ( 12,5 ).
No tocante à Organização Acadêmica, os CEFET/IFET e os Centros Universitários obtiveram Média mais alta que a nacional ( 53,7 e 47,7, respectivamente ).
A Média das Universidades e a Média das Faculdades foram menores do que a nacional ( 46,8 e 45,4, respectivamente ).
Constata -se que existe diferença estatisticamente significativa, ao nível de 95 \%, entre a Média mais alta, 53,7 dos CEFET/IFET, e a dos demais tipos de Organização Acadêmica.
A Média da Modalidade Educação Presencial ( 47,3 ) foi inferior à Média da Educação a Distância ( 48,1 ), e a diferença entre elas não é estatisticamente significativa, ao nível de 95 \%. 6.1.2 Estatísticas Básicas no Componente de Formação Geral Na Tabela 6.3, são apresentadas as Estatísticas Básicas em relação ao componente da prova que avalia a Formação Geral dos estudantes concluintes.
Os alunos de todo o Brasil obtiveram desempenho médio de 42,5.
Quanto à variabilidade, o Desvio padrão das notas dos estudantes do Brasil, como um todo, foi 17,7.
A maior Média foi obtida na região Nordeste ( 43,2 ), e a menor, na região Norte ( 39,2 ).
As demais médias foram : 42,3, na região Sudeste, 42,9, na região Sul e 41,4, na região Centro-Oeste.
Já os maiores valores de Desvio padrão foram obtidos nas regiões Norte, Nordeste e Centro-Oeste ( 18,2 ), e o menor, na região Sudeste ( 16,9 ).
Na região Sul, o Desvio padrão foi 17,8.
A maior nota no Componente de Formação Geral da prova do Enade/2018 foi 94,0, obtida por, pelo menos, um aluno na região Nordeste.
A menor nota Máxima foi obtida na região Norte ( 78,3 ).
A Mediana do Brasil, como um todo, foi 42,1, sendo a menor Mediana encontrada na região Norte ( 39,0 ), e a maior, na região Nordeste ( 42,7 ).
A nota Mínima nessa parte foi zero em todas as regiões, com exceção da região Sudeste, cuja nota Mínima foi 3,6.
Considerando -se as notas, segundo Grande Região, observa -se que não existe diferença estatisticamente significativa entre as Médias das notas do Componente de Formação Geral obtidas nas cinco Grandes Regiões.
No Gráfico 6.2, é possível verificar a avaliação do desempenho dos estudantes do Componente de Formação Geral, a partir do histograma da distribuição das notas correspondentes.
A distribuição é unimodal, com moda em ( 40 ; 50 ], mesmo intervalo ao modal da distribuição de notas da prova, como um todo ( Gráfico 6.1 ).
Nota -se, ainda, que, no Gráfico 6.2, as notas apresentam maior dispersão do que as no Gráfico 6.1 ( distribuição das notas da prova ), confirmado pela comparação dos desvios padrão : 12,5 para a nota da prova, como um todo, e 17,7 para o Componente de Formação Geral.
Na Tabela 6.4, são apresentadas as informações referentes ao desempenho dos concluintes do Componente de Formação Geral, em diferentes agregações : Categoria Administrativa, Organização Acadêmica e Modalidade de Ensino.
Observa -se que existe diferença estatisticamente significativa entre as médias dos tipos de Categoria Administrativa : a maior Média obtida por estudantes de IES Públicas ( 49,5 ) e a menor, pelos de IES Privadas ( 42,0 ).
Nota -se que, também, há diferença estatisticamente significativa entre a maior das médias dos quatro tipos de Organização Acadêmica – CEFET/IFET ( 50,3 ) e as médias das demais, bem como entre as médias dos Centros Universitários e as das Faculdades.
No que se refere à Modalidade de Ensino, a Educação a Distância apresentou maior Média do que a Educação Presencial ( médias 46,3 e 42,3, respectivamente ), verificando -se diferença estatisticamente significativa entre tais médias. 6.1.3 - Estatísticas Básicas do Componente de Conhecimento Específico Na Tabela 6.5, são apresentadas as Estatísticas Básicas referentes ao Componente de Conhecimento Específico da área de Tecnologia em Design de Interiores.
A Média do desempenho dos alunos do Brasil, como um todo, foi 48,9.
A maior Média foi obtida na região Sul ( 49,7 ), e a menor, na região Centro-Oeste ( 47,6 ).
As demais médias foram : 47,7, na região Norte, 49,4, na região Nordeste e 48,4, na região Sudeste.
Quanto à variabilidade das notas, o Desvio padrão do Brasil, como um todo, foi 13,0, sendo o maior Desvio padrão observado na região Norte ( 14,8 ), e o menor, na região Sul ( 12,5 ).
Os demais desvios foram : 13,0, na região Nordeste, 13,1, na região Sudeste e 13,6, na região Centro-Oeste.
A Mediana das notas dos estudantes de todo o Brasil foi 49,7.
A maior Mediana ocorreu na região Sul ( 50,6 ), e a menor, na região Norte ( 47,3 ).
As demais medianas foram : 50,1 na região Nordeste, 49,1, na região Sudeste e 48,8, na região Centro-Oeste.
A nota Máxima do Brasil, como um todo, foi 85,1, sendo obtida por, pelo menos, um aluno na região Centro-Oeste.
As demais notas máximas foram : 74,3, na região Norte, 81,1, na região Nordeste, 81,6, na região Sudeste e 78,1, na região Sul.
A nota Mínima foi zero, obtida nas regiões Sul e Centro-Oeste.
Nas demais regiões, as notas mínimas foram : 6,1, na região Norte, 1,5, na região Nordeste e 1,2, na região Sudeste.
Observa -se que não existe diferença estatisticamente significativa entre as Médias, das notas do Componente de Conhecimento Específico, das cinco Grandes Regiões.
Assim como nos Gráficos 6.1 e 6.2, no Gráfico 6.3, apresentado a seguir, é possível ser feita uma avaliação do desempenho de concluintes de Tecnologia em Design de Interiores, em relação ao Componente de Conhecimento Específico, com um histograma da distribuição das notas correspondentes.
Esta também é uma distribuição unimodal, e o grupo modal é o ( 50 ; 60 ], acima do grupo modal para a Formação Geral e um pouco abaixo do grupo modal da prova, como um todo.
Na Tabela 6.6, é apresentada uma comparação dos resultados em relação à Categoria Administrativa, à Organização Acadêmica e à Modalidade de Ensino, agora levando -se em conta o desempenho de alunos do Componente de Conhecimento Específico da prova. 6.2 ANÁLISE DAS QUESTÕES OBJETIVAS Esta seção apresenta estatísticas selecionadas e histogramas das Questões Objetivas de Formação Geral ( 6.2.1 ) e de Conhecimento Específico ( 6.2.2 ).
São, também, apresentadas e comparadas as médias das subpopulações caracterizadas por Grande Região. 6.2.1 Componente de Formação Geral Na Tabela 6.7, são apresentadas as Estatísticas Básicas relativas às oito questões objetivas do componente da prova que abrange a Formação Geral dos estudantes.
A Média do Brasil foi 47,2.
A menor Média foi encontrada na região Centro-Oeste ( 45,9 ), e a maior, na região Nordeste ( 48,3 ).
As demais médias foram : 46,2, na região Norte, 46,7, na região Sudeste e 47,3, na região Sul.
O Desvio padrão do Brasil foi 22,6, sendo o maior Desvio CEFET/IFET obtiveram nota Máxima 80,9.
As medianas foram : 49,1, nas Universidades ; 50,3, nos Centros Universitários ; 46,6, nas Faculdades, a menor delas ; e 56,1 nos CEFET/IFET, a maior.
A nota Mínima foi zero, em Centros Universitários.
Nas demais Organizações Acadêmicas, as notas mínimas foram : 1,2, nas Universidades ; 3,6, nas Faculdades e 26,6, nos CEFET/IFET.
Observa -se que existe diferença estatisticamente significativa, ao nível de 95 \%, no Componente de Conhecimento Específico entre as notas dos CEFET/IFET e as das Universidades, dos Centros Universitários e das Faculdades.
Quanto à Categoria Administrativa, observa -se um comportamento semelhante àquele da parte de Formação Geral e à prova, como um todo, ou seja, existe diferença estatisticamente significativa entre as médias das IES Públicas ( 54,3 ) e as das IES Privadas ( 48,5 ).
Neste caso, também, a maior Média foi obtida por alunos de IES Públicas de ensino.
Quanto ao comportamento das médias no que toca à Modalidade de Ensino, não há diferença estatisticamente significativa entre a Média da Educação Presencial ( 49,0 ) e a da Educação a Distância ( 48,7 ). 6.2 ANÁLISE DAS QUESTÕES OBJETIVAS Esta seção apresenta estatísticas selecionadas e histogramas das Questões Objetivas de Formação Geral ( 6.2.1 ) e de Conhecimento Específico ( 6.2.2 ).
São, também, apresentadas e comparadas as médias das subpopulações caracterizadas por Grande Região. 6.2.1 Componente de Formação Geral Na Tabela 6.7, são apresentadas as Estatísticas Básicas relativas às oito questões objetivas do componente da prova que abrange a Formação Geral dos estudantes.
A Média do Brasil foi 47,2.
A menor Média foi encontrada na região Centro-Oeste ( 45,9 ), e a maior, na região Nordeste ( 48,3 ).
As demais médias foram : 46,2, na região Norte, 46,7, na região Sudeste e 47,3, na região Sul.
O Desvio padrão do Brasil foi 22,6, sendo o maior Desvio padrão encontrado na região Norte ( 23,6 ), e o menor, nas regiões Sudeste e Centro-Oeste ( 22,0 ).
Os demais desvios foram : 23,3, na região Nordeste e 22,7, na região Sul.
As medianas do Brasil, como um todo, e de todas as cinco regiões foram 50,0.
A nota Máxima 100,0 foi alcançada em quatro das Grandes Regiões.
A exceção foi a região Norte, com nota Máxima 87,5.
As notas mínimas ( 0,0 ) foram iguais para todas as regiões.
Na Tabela 6.8, são apresentados o Índice de Facilidade e o Índice de Discriminação ( Ponto-Bisserial ) para \textbf{cada} uma das questões objetivas do Componente de Formação Geral.
Quanto ao Índice de Facilidade, foram usadas as seguintes cores para diferenciar o nível de dificuldade da questão :  Azul para as questões classificadas com índice Muito fácil ( \&gt ; =0,86 ), verde para as questões classificadas com índice Fácil ( 0,61 a 0,85 ), amarelo para as questões classificadas com Médio ( 0,41 a 0,60 ), vermelho para as questões classificadas com Difícil ( 0,16 a 0,40 ) e roxo para as questões classificadas com Muito difícil ( \&lt ; =0,15 ).
Já quanto ao Índice de Discriminação, foram usadas as seguintes cores para qualificar a questão :  As questões classificadas com índice Fraco receberam a cor vermelha ( \&lt ; =0,19 ), as classificadas com Médio receberam a cor amarela ( 0,20 a 0,29 ), as classificadas com Bom receberam a cor verde ( 0,30 a 0,39 ) e as classificadas com Muito bom ( \&gt ; =0,40 ) receberam a cor azul.
As questões objetivas do Componente de Formação Geral, segundo o Índice de Facilidade, foram assim avaliadas : das oito questões, nenhuma teve o Índice de Facilidade classificado como Muito fácil, e uma questão foi tida como Fácil, por ter índice de acertos 0,63.
Seis questões foram consideradas com índice de dificuldade Médio, situando -se no intervalo entre 0,43 e 0,57 do Índice de Facilidade, ou seja, houve entre 43,0 \% e 57,0 \% de acertos, enquanto uma questão foi classificada na categoria Difícil, com Índice de Facilidade 0,25. padrão encontrado na região Norte ( 23,6 ), e o menor, nas regiões Sudeste e Centro-Oeste ( 22,0 ).
Os demais desvios foram : 23,3, na região Nordeste e 22,7, na região Sul.nAs medianas do Brasil, como um todo, e de todas as cinco regiões foram 50,0.
A nota Máxima 100,0 foi alcançada em quatro das Grandes Regiões.
A exceção foi a região Norte, com nota Máxima 87,5.
As notas mínimas ( 0,0 ) foram iguais para todas as regiões.nnNa Tabela 6.8, são apresentados o Índice de Facilidade e o Índice de Discriminação ( Ponto-Bisserial ) para \textbf{cada} uma das questões objetivas do Componente de Formação Geral.
Quanto ao Índice de Facilidade, foram usadas as seguintes cores para diferenciar o nível de dificuldade da questão:n Azul para as questões classificadas com índice Muito fácil ( \&gt ; =0,86 ), verde para as questões classificadas com índice Fácil ( 0,61 a 0,85 ), amarelo para as questões classificadas com Médio ( 0,41 a 0,60 ), vermelho para as questões classificadas com Difícil ( 0,16 a 0,40 ) e roxo para as questões classificadas com Muito difícil ( \&lt ; =0,15).nJá quanto ao Índice de Discriminação, foram usadas as seguintes cores para qualificar a questão:n As questões classificadas com índice Fraco receberam a cor vermelha ( \&lt ; =0,19 ), as classificadas com Médio receberam a cor amarela ( 0,20 a 0,29 ), as classificadas com Bom receberam a cor verde ( 0,30 a 0,39 ) e as classificadas com Muito bom ( \&gt ; =0,40 ) receberam a cor azul.nAs questões objetivas do Componente de Formação Geral, segundo o Índice de Facilidade, foram assim avaliadas : das oito questões, nenhuma teve o Índice de Facilidade classificado como Muito fácil, e uma questão foi tida como Fácil, por ter índice de acertos 0,63.
Seis questões foram consideradas com índice de dificuldade Médio, situando -se no intervalo entre 0,43 e 0,57 do Índice de Facilidade, ou seja, houve entre 43,0 \% e 57,0 \% de acertos, enquanto uma questão foi classificada na categoria Difícil, com Índice de Facilidade 0,25.
Nenhuma questão apresentou menos de 15 \% de acertos, razão pela qual não houve questão classificada como Muito difícil.
O Índice de Facilidade variou de 0,25 a 0,63, e o de Discriminação, de 0,42 a 0,49.
Na Tabela 6.9, é apresentada a distribuição das questões, levando -se em conta, simultaneamente, a classificação dos dois índices.
Todas as questões que tiveram Índice de Discriminação Muito bom figuraram entre três níveis de dificuldade Fácil, Médio e Difícil : uma classificada na categoria Fácil ( questão 2 ) do Índice de Facilidade, seis, na categoria Médio ( questões 1, 3, 5, 6, 7, 8 ), e uma, na categoria Difícil ( questão 4 ).
Em particular, as questões 4 e 5 apresentaram o maior poder discriminatório, com índice 0,49, sendo a 4 considerada Difícil, e a 5, de grau Médio em termos de facilidade, com uma proporção de acertos de, respectivamente, 0,25 e 0,44.
O máximo de acertos foi alcançado pela questão 2 com um Índice de Facilidade de 0,63.
As demais questões, de grau Médio, apresentaram os seguintes Índice de Facilidade : questão 1, 0,46 ; questão 3, 0,48 ; questão 6, 0,43 ; questão 7, 0,51 ; e questão 8, 0,57.
Já os índices de Discriminação foram Muito bom para todas as 8 questões, com valores variando de 0,42 a 0,49.
A questão de número 4 apresentou um Índice de Facilidade de 0,25 ( Difícil ), ou seja, um quarto dos estudantes conseguiu resolvê -la, dentro do universo de participantes.
Seu Índice de Discriminação foi Muito Bom ( 0,49 ).
No Gráfico 6.4, para exemplificar, analisa -se o comportamento da questão de número 4 de Formação Geral.
Trata -se de uma questão considerada Difícil, em relação à facilidade e a que obteve o maior Índice de Discriminação dessa parte da prova.
Neste gráfico, \textbf{cada} uma das cinco curvas representa o percentual de respostas em determinada alternativa da questão, em função do número de acertos dos estudantes nessa parte da prova ( Formação Geral/Múltipla Escolha ), antes de possíveis eliminações pelo critério do Ponto-Bisserial.
A curva em verde corresponde à alternativa B, a alternativa correta para essa questão.
Observa -se que entre os estudantes com menor número de acertos, nessa parte do exame, a situação mais frequente foi a escolha da alternativa C ( em preto ), incorreta.
Por exemplo, dentre os estudantes que acertaram três questões, 38,4 \% escolheram a alternativa C ( em preto ), 22,1 \% escolheram a alternativa E ( em vermelho ), 17,0 \% escolheram a alternativa D ( em laranja ), 13,2 \% escolheram a alternativa B ( em verde ), e 8,8 \%, a A ( azul ).
Dentre os que acertaram três respostas entre as questões de múltipla escolha de Formação Geral, 0,2 \% deixou a questão em branco e 0,2 \% marcou mais de uma alternativa, invalidando a questão.
À medida que o número de acertos aumenta, indicando desempenho melhor nessa parte da prova, aumenta, concomitantemente, a proporção de estudantes que selecionaram a alternativa correta B, atingindo 100 \% para os estudantes com oito acertos.
Essa análise permite verificar como a questão discriminou os grupos de desempenho, justificando -se o alto índice obtido na questão.
Cumpre notar que não é possível inferir deste gráfico nem o índice de facilidade ( que seria uma média da proporção ponderada pela quantidade de alunos com \textbf{cada} uma das notas ), nem o índice de discriminação Ponto-Bisserial, por razão equivalente.
No caso extremo, no qual a grande concentração dos acertos dos alunos fosse abaixo de quatro, o índice de facilidade seria, obrigatoriamente, abaixo de cerca de 20 \% ( neste exemplo ).
Caso a concentração fosse em seis acertos ou mais, o índice seria, obrigatoriamente, acima de 50 \%.
Os gráficos relativos às demais questões de Formação Geral constam do Anexo I.
No Gráfico 6.4, para exemplificar, analisa -se o comportamento da questão de número 4 de Formação Geral.
Trata -se de uma questão considerada Difícil, em relação à facilidade e a que obteve o maior Índice de Discriminação dessa parte da prova.
Neste gráfico, \textbf{cada} uma das cinco curvas representa o percentual de respostas em determinada alternativa da questão, em função do número de acertos dos estudantes nessa parte da prova ( Formação Geral/Múltipla Escolha ), antes de possíveis eliminações pelo critério do Ponto-Bisserial.
A curva em verde corresponde à alternativa B, a alternativa correta para essa questão.
Observa -se que entre os estudantes com menor número de acertos, nessa parte do exame, a situação mais frequente foi a escolha da alternativa C ( em preto ), incorreta.
Por exemplo, dentre os estudantes que acertaram três questões, 38,4 \% escolheram a alternativa C ( em preto ), 22,1 \% escolheram a alternativa E ( em vermelho ), 17,0 \% escolheram a alternativa D ( em laranja ), 13,2 \% escolheram a alternativa B ( em verde ), e 8,8 \%, a A ( azul ).
Dentre os que acertaram três respostas entre as questões de múltipla escolha de Formação Geral, 0,2 \% deixou a questão em branco e 0,2 \% marcou mais de uma alternativa, invalidando a questão.
À medida que o número de acertos aumenta, indicando desempenho melhor nessa parte da prova, aumenta, concomitantemente, a proporção de estudantes que selecionaram a alternativa correta B, atingindo 100 \% para os estudantes com oito acertos.
Essa análise permite verificar como a questão discriminou os grupos de desempenho, justificando -se o alto índice obtido na questão.
Cumpre notar que não é possível inferir deste gráfico nem o índice de facilidade ( que seria uma média da proporção ponderada pela quantidade de alunos com \textbf{cada} uma das notas ), nem o índice de discriminação Ponto-Bisserial, por razão equivalente.
No caso extremo, no qual a grande concentração dos acertos dos alunos fosse abaixo de quatro, o índice de facilidade seria, obrigatoriamente, abaixo de cerca de 20 \% ( neste exemplo ).
Caso a concentração fosse em seis acertos ou mais, o índice seria, obrigatoriamente, acima de 50 \%.
Os gráficos relativos às demais questões de Formação Geral constam do Anexo I. 6.2.2 Componente de Conhecimento Específico Na Tabela 6.10, são apresentadas as Estatísticas Básicas em relação às questões objetivas do Componente de Conhecimento Específico da prova de Tecnologia em Design de Interiores por Grande Região.
Três questões objetivas do Componente de Conhecimentos Específicos não foram usadas no cômputo das notas por terem sido descartadas pelo critério Ponto-Bisserial.
Assim, as notas foram calculadas com base em 24 das 27 questões objetivas de Conhecimentos Específicos.
A Média do Brasil desse componente foi 52,6.
A menor Média foi observada na região Centro-Oeste ( 51,8 ), e a maior, nas regiões Nordeste e Sul ( 53,1 ).
O Desvio padrão de todo o Brasil foi 14,0, sendo o menor Desvio padrão encontrado na região Sul ( 13,5 ), e o maior, na região Norte ( 15,6 ).
A Mediana de todo o Brasil foi 54,2, o mesmo valor da Mediana encontrada em todas as regiões.
A nota Máxima da prova foi 91,7, obtida nas questões objetivas do Componente de Conhecimento Específico, por, pelo menos, um aluno da região Sudeste.
A nota Mínima foi zero, obtida em todas as regiões, exceto na Norte ( 4,2 ).
Na Tabela 6.11, são apresentados os Índices de Facilidade e Discriminação ( Ponto-Bisserial ) das questões objetivas do Componente de Conhecimento Específico, para os estudantes de Tecnologia em Design de Interiores.
Para facilitar a diferenciação das questões, serão usadas as mesmas cores da Tabela 6.8 para as diferentes classificações dos Índices de Facilidade e de Discriminação.
As análises dizem respeito às 27 questões válidas.
A partir do Índice de Facilidade obtido, pode -se concluir que cerca de um terço das questões objetivas da prova foi considerado, pelo menos, Difícil : das 27 questões válidas, 8 foram classificadas como Difícil, e uma, como Muito difícil.
Duas questões foram classificadas como Muito fácil, quatro, como Fácil, e outras 12 consideradas Médio.
Já quanto aos índices de discriminação das questões objetivas do Componente de Conhecimento Específico da prova, obtém -se como resultado a seguinte classificação : 14 das 27 questões foram consideradas boas.
Assim, para um pouco mais da metade das questões – 14 em 27 – o Índice de Discriminação foi Bom.
Dentre as demais, dez delas foram classificadas como Médio, e outras três, como Fraco, sendo 13, por conseguinte, a quantidade de questões nos dois patamares mais baixos de discriminação.
Constata -se, assim, que a prova – no que se refere ao Componente de Conhecimento Específico – possuía média capacidade de discriminar dentre aqueles que dominam ou não o conteúdo.
Na Tabela 6.12, é apresentada a distribuição das questões levando -se em conta, simultaneamente, a classificação dos dois índices.
Dentre as questões que alcançaram os maiores índices de discriminação, Bom, duas foram classificadas, quanto ao índice de Facilidade, como Difícil : as questões 16 e 25, com, respectivamente, 38 \% e 36 \% dos estudantes marcando a opção correta.
As outras 12 questões que alcançaram Índices de Discriminação Bom distribuíram -se da seguinte forma no que se refere ao Índice de Facilidade : sete classificadas como Médio ; quatro, como Fácil ; e uma, como Muito fácil.
Dessas 12 classificadas como Bom, a que alcançou o maior valor foi a questão 32, classificada como Fácil e com Índice de Discriminação 0,39.
A questão de número 21 foi a mais difícil dentre as 27 questões específicas válidas, com baixo Índice de Facilidade, apenas 14 \% de acertos.
Essa questão apresentou poder discriminatório médio, 0,21.
As outras nove questões com Índice de Discriminação Médio, Na Tabela 6.12, é apresentada a distribuição das questões levando -se em conta, simultaneamente, a classificação dos dois índices.
Dentre as questões que alcançaram os maiores índices de discriminação, Bom, duas foram classificadas, quanto ao índice de Facilidade, como Difícil : as questões 16 e 25, com, respectivamente, 38 \% e 36 \% dos estudantes marcando a opção correta.
As outras 12 questões que alcançaram Índices de Discriminação Bom distribuíram -se da seguinte forma no que se refere ao Índice de Facilidade : sete classificadas como Médio ; quatro, como Fácil ; e uma, como Muito fácil.
Dessas 12 classificadas como Bom, a que alcançou o maior valor foi a questão 32, classificada como Fácil e com Índice de Discriminação 0,39.
A questão de número 21 foi a mais difícil dentre as 27 questões específicas válidas, com baixo Índice de Facilidade, apenas 14 \% de acertos.
Essa questão apresentou poder discriminatório médio, 0,21.
As outras nove questões com Índice de Discriminação Médio, distribuíram -se da seguinte forma quanto ao Índice de Facilidade : quatro difíceis, quatro médias e uma muito fácil.
Três questões apresentaram Índice de Discriminação Fraco.
Tais questões foram, portanto, pelo critério Ponto-Bisserial, consideradas inadequadas.
Por isso, as questões 18, 26 e 35 foram eliminadas do cômputo da nota final.
A título de exemplo das análises do comportamento das questões objetivas, no Gráfico 6.5, analisa -se a questão 32 do Componente de Conhecimento Específico.
Essa questão foi considerada, pelas respostas dos estudantes avaliados na prova, como uma questão com nível de facilidade Fácil, 0,83, ou seja, 83 \% dos estudantes assinalaram acertadamente a opção E, correspondente ao gabarito.
Como já comentado, seu Índice de Discriminação foi igual a 0,39, classificado como Bom, o maior valor de discriminação.
Neste gráfico, \textbf{cada} uma das cinco curvas representa o percentual de respostas em determinada alternativa da questão 32, em função do número de acertos dos estudantes nessa parte da prova, antes de possíveis eliminações de questões pelo critério Ponto-Bisserial.
A alternativa correta E, representada no gráfico pela curva em vermelho, foi escolhida em maiores proporções pelos alunos com desempenho melhor nessa parte da prova.
Já as alternativas incorretas, também denominadas distratores, foram selecionadas, principalmente, por aqueles com notas mais baixas.
Observa -se que a soma não é 100 \%, por causa das questões não respondidas ou com mais de uma opção marcada.
Aqueles com nota zero, na sua quase totalidade, deixaram esta questão em branco ou marcaram mais de uma alternativa, comportamento considerado inválido.
A proporção de alunos que selecionaram a resposta correta E aumenta gradativamente, chegando a atingir 100 \% para 20 acertos ou mais, enquanto a proporção dos que escolheram alternativas incorretas decai, a partir de cinco acertos, em função do número de acertos nessa parte da prova.
Os gráficos relativos às demais questões do Conhecimento Específico constam do Anexo I. 6.3 ANÁLISE DAS QUESTÕES DISCURSIVAS Esta seção apresenta estatísticas selecionadas e histogramas das Questões Discursivas de Formação Geral ( 6.3.1 ) e Conhecimento Específico ( 6.3.2 ).
São também apresentadas e comparadas as médias de subpopulações, caracterizadas por Grande Região. 6.3.1 Componente de Formação Geral As análises dos resultados de desempenho dos estudantes de Tecnologia em Design de Interiores, nas duas questões discursivas relativas à Formação Geral, encontram -se na Tabela 6.13 e no Gráfico 6.6.
Na Tabela 6.13, observa -se que a nota Média nesse conjunto de questões foi abaixo da obtida nas objetivas.
Os estudantes de todo Brasil obtiveram, em Formação Geral, Média 47,2, nas questões objetivas e 35,6, nas questões discursivas.
Pode -se notar, também, que o Desvio padrão nesse conjunto de questões foi menor do que o obtido nas objetivas : 22,6, nas questões objetivas e 20,1, nas questões discursivas.
A maior Média foi obtida na região Sul ( 36,4 ), e a menor, na região Norte ( 28,9 ).
A Mediana de todo o Brasil, neste componente, foi 37,5, a mesma sendo obtida nas regiões Sul e Centro-Oeste.
Na região Norte, a Mediana foi menor ( 31,0 ), e, nas demais regiões, foi maior que o valor para o Brasil : 38,0, nas regiões Nordeste e Sudeste.
A nota Máxima ( 89,5 ) foi obtida na região Sudeste, sendo 77,0, a Máxima na região Norte ; 86,5, na Nordeste ; 84,5, na Sul ; e 88,0, na região Centro-Oeste.
A nota Mínima ( 0,0 ) foi a mesma em todas as regiões do Brasil.
No Gráfico 6.6, \textbf{está} representada a distribuição das notas nas questões discursivas do Componente de Formação Geral.
É distribuição bimodal, com as modas ocorrendo nos intervalos ( 30 ; 40 ] e no intervalo ( 40 ; 50 ], ambas com frequência 19,9 \%.
Destaca -se, também, o intervalo [ 0 ; 10 ], com distribuição em \textbf{torno} de 13,7 \% do total de notas, sendo que no intervalo [ 0 ; 10 ], inclui -se, além da nota zero, a frequência de alunos que deixaram esse tipo de questão em branco.
Na sequência, os resultados verificados para \textbf{cada} uma das questões discursivas de Formação Geral serão apresentados, estabelecendo -se relações com os conteúdos abordados em \textbf{cada} uma delas.
Os comentários da Banca de docentes corretores, a respeito do observado na correção das respostas dos estudantes, suas impressões e conclusões serão apresentados junto à análise de \textbf{cada} questão.
Cumpre esclarecer que, tendo -se em vista que as questões discursivas de Formação Geral são padronizadas, ou seja, constam de todas as provas, os comentários da Banca são os mesmos para todas as carreiras acadêmicas, sendo direcionados a todos os estudantes que participaram do Enade/2018.
A seguir, serão analisados os desempenhos da Área de Tecnologia em Design de Interiores nas duas questões discursivas de Formação Geral do Enade/2018, comparando -se os resultados obtidos com comentários para \textbf{cada} questão. 6.3.1.1 Análise de Conteúdo da Questão Discursiva 1 do Componente de Formação Geral Os dados de Tecnologia em Design de Interiores, obtidos a partir das respostas à questão 1, encontram -se na Tabela 6.14 e no Gráfico 6.7.
Nessa questão, os alunos de todo o Brasil tiveram Média 29,9, inferior à Média obtida na questão discursiva 2 ( média 32,5 ).
A maior Média para a questão 1 foi obtida na região Nordeste ( 31,2 ), e a menor, na região Norte ( 23,2 ).
Quanto à variabilidade das notas, o Desvio padrão de todo o Brasil foi 25,5, o mesmo da região Sul.
O menor Desvio padrão foi obtido na região Sudeste ( 24,5 ), e o maior Desvio padrão foi obtido na região Nordeste ( 26,6 ).
A mediana da região Norte foi 10,0, e nas outras quatro regiões e no Brasil, como um todo foi 30,0.
As notas máximas da questão discursiva 1 do Brasil e de quatro regiões foi 100,0, com apenas a região Centro-Oeste apresentando valor diferente : 90,0.
Já a nota Mínima da questão discursiva 1 foi, sem exceção, zero para todas as regiões do Brasil.
No Gráfico 6.7, mostra -se a distribuição das notas na questão discursiva 1 do Componente de Formação Geral.
Observa -se que a maior frequência corresponde ao intervalo ( 30 ; 40 ], correspondendo a 16,2 \%, seguido de perto pela dos alunos que deixaram a resposta a essa questão em branco, 15,8 \% dos participantes. 6.3.1.2 Comentários sobre a correção de Conteúdo das respostas à Questão Discursiva 1 A Questão 1 exigia do aluno um posicionamento sobre conteúdo humanista e crítico.
O comando da questão solicitava a redação de um texto que formulasse e articulasse argumentos consistentes envolvendo temas previstos na Portaria nº 444, de 30 de maio de 2018, que dispõe sobre o componente de Formação geral do Enade 2018, tais como : solidariedade/violência, tolerância/intolerância, inclusão/exclusão, sexualidade, relações de gênero e relações étnico-raciais ( art. 7º, item VIII ), bem como ética, democracia e cidadania ( art. 7º, item I ).
O enunciado da questão buscava o estabelecimento da relação de três dimensões – os Direitos Humanos, a Democracia e as Fake News.
Para isso, recorreu a quatro suportes : a reprodução de uma foto e três fragmentos adaptados de textos que foram disponibilizados em páginas da Internet ( da Anistia Internacional, do Jornal El País-Brasil e do Labic-Laboratório de estudos sobre Imagem e Cibercultura ).
Os fragmentos traziam dados sobre o aumento do número de assassinatos de grupos minoritários, as origens e o trabalho da vereadora Marielle, sua luta pelos Direitos Humanos, seu assassinato, e a produção de fake news.
O comando solicitava um texto abordando dois tópicos : o tensionamento entre a defesa dos Direitos Humanos e a produção de notícias falsas, tendo como base a vereadora Marielle – a defesa dos Direitos Humanos realizada por ela e as notícias falsas divulgadas após a sua morte.
As fake news também \textbf{estavam} presentes no segundo tópico de análise : os prejuízos que provocam para a sociedade democrática.
Por um lado, os textos selecionados para constituir o enunciado da questão continham apoio suficiente para formulação e articulação de argumentos para discorrer sobre o tensionamento entre a atuação da Marielle em defesa dos Direitos Humanos e fake news.
Por outro, havia ênfase nos aspectos biográficos e na trajetória da vereadora Marielle, o que levou a respostas pouco focadas no comando, que discutiam a vida e a morte da vereadora, sua importância como um ícone dos Direitos Humanos, opiniões em \textbf{torno} do seu assassinato, até mesmo sobre a responsabilidade pelo crime, muitos textos com a presença/defesa de fake news ou produção de novas.
O foco sobre Marielle também ensejou respostas basicamente centradas na questão da violência.
Foram frequentes as interpretações com viés político-partidário, não previstas no padrão de resposta.
O item do comando que solicitava a discussão sobre democracia e fake news encontrou menos apoio nos textos.
Apesar de se esperar que graduandos tivessem maturidade, inclusive política, para demonstrar alguma compreensão dos aspectos principais que envolvem o convívio democrático, esse aspecto foi pouco discutido e, quando presente, quase sempre foi apenas mencionado, numa quase reprodução do comando.
A falta de entendimento do que é o conceito de Direitos Humanos também foi percebida nas respostas, seja para recomendá -lo como importante para a sociedade, seja para criticá -lo como algo que defende bandidos.
Foi marcante a ideia que vincula Direitos Humanos a direitos apenas das minorias.
Não foi incomum encontrar respostas que personificavam os Direitos Humanos, tratando -os como sujeito de alguma ação que deveria ser implementada na sociedade em geral ou em casos específicos.
As fake news foram citadas muitas vezes, mas pouco explicadas, conforme previsto no padrão de resposta.
O argumento utilizado geralmente vinculava as consequências de fake news para a imagem da Marielle, ou seja, centrava no fato em si, já indicado no enunciado, sem qualquer análise ou aprofundamento próprio.
É importante salientar que muitas respostas, provavelmente de grupos vinculados a cursos específicos, continham argumentações recorrentes : os Direitos Humanos como algo vinculado à Constituição e a tratados dos quais o Brasil é signatário ; as fake news como vereadora Marielle Franco, sua luta pelos Direitos Humanos, seu assassinato, e a produção de fake news.
O comando solicitava um texto abordando dois tópicos : o tensionamento entre a defesa dos Direitos Humanos e a produção de notícias falsas, tendo como base a vereadora Marielle – a defesa dos Direitos Humanos realizada por ela e as notícias falsas divulgadas após a sua morte.
As fake news também \textbf{estavam} presentes no segundo tópico de análise : os prejuízos que provocam para a sociedade democrática.
Por um lado, os textos selecionados para constituir o enunciado da questão continham apoio suficiente para formulação e articulação de argumentos para discorrer sobre o tensionamento entre a atuação da Marielle em defesa dos Direitos Humanos e fake news.
Por outro, havia ênfase nos aspectos biográficos e na trajetória da vereadora Marielle, o que levou a respostas pouco focadas no comando, que discutiam a vida e a morte da vereadora, sua importância como um ícone dos Direitos Humanos, opiniões em \textbf{torno} do seu assassinato, até mesmo sobre a responsabilidade pelo crime, muitos textos com a presença/defesa de fake news ou produção de novas.
O foco sobre Marielle também ensejou respostas basicamente centradas na questão da violência.
Foram frequentes as interpretações com viés político-partidário, não previstas no padrão de resposta.
O item do comando que solicitava a discussão sobre democracia e fake news encontrou menos apoio nos textos.
Apesar de se esperar que graduandos tivessem maturidade, inclusive política, para demonstrar alguma compreensão dos aspectos principais que envolvem o convívio democrático, esse aspecto foi pouco discutido e, quando presente, quase sempre foi apenas mencionado, numa quase reprodução do comando.
A falta de entendimento do que é o conceito de Direitos Humanos também foi percebida nas respostas, seja para recomendá -lo como importante para a sociedade, seja para criticá -lo como algo que defende bandidos.
Foi marcante a ideia que vincula Direitos Humanos a direitos apenas das minorias.
Não foi incomum encontrar respostas que personificavam os Direitos Humanos, tratando -os como sujeito de alguma ação que deveria ser implementada na sociedade em geral ou em casos específicos.
As fake news foram citadas muitas vezes, mas pouco explicadas, conforme previsto no padrão de resposta.
O argumento utilizado geralmente vinculava as consequências de fake news para a imagem da Marielle, ou seja, centrava no fato em si, já indicado no enunciado, sem qualquer análise ou aprofundamento próprio.
É importante salientar que muitas respostas, provavelmente de grupos vinculados a cursos específicos, continham argumentações recorrentes : os Direitos Humanos como algo vinculado à Constituição e a tratados dos quais o Brasil é signatário ; as fake news como Como já apontado, houve baixa quantidade de respostas contemplando o segundo tópico do comando – prejuízos da produção de notícias falsas para a sociedade democrática ; fato motivado, possivelmente, pela falta de um elemento que ressaltasse esse ponto nos textos motivadores.
De modo geral, os estudantes não conseguiram recorrer a conhecimentos tácitos sobre democracia para construir argumentos próprios.
As respostas que abordavam o tema, geralmente \textbf{estavam} vinculadas à ideia do direito à verdade, do direito à integridade da pessoa ou dos direitos gerais garantidos pela Constituição.
A relação da questão com o momento político brasileiro foi indicada em um número expressivo de respostas.
Por vezes, apenas questionando a temática da questão – caso Marielle – em uma prova para avaliar o ensino superior no Brasil ou a importância dos Direitos Humanos na atual conjuntura política do país, levando à anulação da resposta.
Houve também interpretações das consequências do uso das fake news no processo eleitoral de 2018, influenciando a escolha consciente e livre das pessoas.
As respostas demonstraram atenção político social frente a esses temas, porém com pouco domínio conceitual do campo de conhecimento sobre Direitos Humanos, fake news e democracia.
Em termos da linguagem adotada pelos alunos, há dois grandes grupos de respostas.
No primeiro grupo, os que responderam de forma clara, com autonomia de pensamento e explicitação de justificativas para os aspectos solicitados no enunciado, demonstrando conhecimento sobre o que significam Direitos Humanos, fake news e democracia, além da articulação entre essas questões.
Por outro lado, um segundo grupo dos que não conseguiram ir além do que foi citado no enunciado, apenas repetindo -o, muitas vezes desviando o foco da questão para temas correlatos, como a vida de Marielle ou a violência.
Nesse segundo grupo, mais facilmente foram encontradas respostas desconexas e mal articuladas, com textos confusos que explicitam o não conhecimento sobre o tema e a falta de adequação à linguagem, necessária a um concluinte do ensino superior.
As abordagens dos Direitos Humanos mais recorrentes nas respostas corretas foram : a referência da Marielle como representante/ícone dos DH ; e os DH como direitos previstos na Constituição e em tratados internacionais dos quais o Brasil é signatário.
No entanto, também foram encontradas críticas aos Direitos Humanos, tais como : DH vistos como “ direitos de bandidos ” ; a defesa pouco embasada de frases de efeito como “ direitos humanos para humanos direitos ”.
Também chamou atenção o tratamento dos DH como um personagem, ou um órgão público, evidenciando desconhecimento sobre o tema.
Em relação às fake news, dois argumentos previstos pelo padrão de respostas foram os mais presentes : a noção de que o mundo globalizado e a facilidade de reprodução de informações pelas mídias atuais levam à criação de mentiras que são compartilhadas e que geram consequências ; a Marielle como alvo de fake news que tentaram desconstruir sua luta e seu legado.
Também foi frequente a ideia de que a produção de informações falsas é realizada por uma imprensa tendenciosa com intenção de manipular a informação.
No item relacionado a fake news e democracia, as respostas corretas afirmavam que as fake news inviabilizam o acesso à verdade e à garantia da integridade de um Estado Democrático de Direito.
Também foi abordada a influência das fake news nos resultados das eleições dos últimos anos, nos Estados Unidos e no Brasil, gerando consequências para a vida das pessoas.
Por fim, um outro aspecto abordado nas respostas a este tópico foi a polarização da sociedade em termos políticos, com a afirmação de que as fake news dificultam o debate, estimulam posições extremas e rivalidades prejudiciais à convivência em uma democracia plena.
A notas mais fracas ( de zero a 35 ) foram cerca de 20 \% das respostas corrigidas – as que não \textbf{estavam} em branco nem foram desconsideradas ou anuladas.
Foram respostas que se restringiam a cópia ou paráfrases do enunciado ; as que tinham foco em desdobramentos políticos e policiais do caso Marielle ; as que tinham argumentação confusa e/ou incompleta.
A grande maioria das notas, em \textbf{torno} de 65 \%, foi mediana, entre 40 a 70.
Nesse grupo ficaram as respostas que continham alguma abordagem adequada sobre Direitos Humanos, sobre as atividades da vereadora Marielle e sobre as fake news divulgadas após sua morte.
As notas mais altas deste grupo foram para respostas que incluíram a questão do tensionamento, quase sempre, em um texto coerente e consistente.
Muitas vezes, a nota mediana foi consequência do baixo índice de respostas à segunda parte do comando, que solicitava a discussão sobre os prejuízos das fake news para a democracia.
As melhores notas só foram obtidas por aqueles que incluíram a discussão sobre os prejuízos das fake news para a democracia.
Apenas 15 \% das respostas corrigidas receberam avaliação igual ou superior a 75.
A deficiência principal detectada pelos corretores foi a dificuldade de grande parte dos alunos para interpretar o enunciado da questão e para se expressar em textos consistentes.
De modo geral, os temas foram tratados de forma fragmentada e pouco analítica.
Além disso, chamou à atenção a predominância do lugar comum na construção argumentativa das respostas, a superficialidade do debate em \textbf{torno} da democracia e o desconhecimento de conceitos como Direitos Humanos, esses confundidos com instituições e até pessoas. 6.3.1.3 Análise de Conteúdo da Questão Discursiva 2 do Componente de Formação Geral A Tabela 6.15 mostra que o desempenho médio dos estudantes na questão discursiva 2 ( média 32,5 ) foi superior ao obtido na questão discursiva 1 ( média 29,9 ).
Nessa questão, a região de maior média foi a Sul ( 33,7 ), enquanto a de menor Média foi a região Norte ( 26,5 ).
Quanto à variabilidade das notas, o Desvio padrão de todo o Brasil foi 23,6, inferior ao obtido na questão discursiva 1 ( 25,5 ).
O maior desvio nessa questão foi obtido na região Nordeste ( 24,3 ), enquanto o menor foi obtido na região Norte ( 21,3 ).
A Mediana de todo o Brasil foi 30,0, a mesma das regiões Sudeste e Centro-Oeste.
Nas demais regiões, as medianas foram 35,0, nas regiões Nordeste e Sul, e 25,0, na Região Norte.
A nota máxima no Brasil foi 100, 0, valor que se repetiu nas regiões Nordeste, Sudeste e Sul.
Nas demais regiões, as notas máximas foram 85,0, na Norte e 90,0, na Centro-Oeste.
A nota mínima ( 0,0 ) foi a mesma em todas as regiões, sem exceção.
No Gráfico 6.8, mostra -se a distribuição das notas na questão discursiva 2 do Componente de Formação Geral.
Observa -se que a maior frequência ( 21,5 \% ) corresponde ao intervalo ( 30 ; 40 ].
A segunda maior frequência é no intervalo ( 40 ; 50 ), com 14,5 \%.
A terceira maior frequência ( 13,6 \% ) corresponde aos alunos que deixaram a questão em branco, praticamente igual à do intervalo ( 50 ; 60 ], 13,5 \%.
Nota -se, ainda, que as notas ficaram um pouco menos dispersas, em comparação às da questão discursiva de número 1, o que pode ser constatado, também, pela comparação do desvio padrão das notas da questão discursiva 2 ( 23,6 ) e o das notas da questão discursiva 1 ( 25,5 ). 6.3.1.4 Comentários sobre a correção de Conteúdo das respostas à Questão Discursiva 2 Os textos e o comando da Questão 2 de Formação Geral são perfeitamente adequados ao que se espera avaliar em Formação Geral, ao tratar das funções sociais dos museus, contextualizando a temática sob o ponto de vista da memória e da perspectiva de futuro.
Na questão, \textbf{estavam} presentes, principalmente, o conteúdo humanista e crítico, a interpretação e a produção de textos sobre temas como cultura e arte, educação, ciência e democracia, previstos na Portaria que dispõe sobre o componente de Formação Geral do Enade 2018.
Os conhecimentos exigidos para responder ao comando, explicitados no padrão de resposta, vão além do senso comum, por se exigir a compreensão de conceitos como “ função social ” ( de instituições sociais ), “ memória ” ( no sentido de memória social ), “ perspectiva de futuro ” ( no contexto sócio-histórico-cultural ) e “ sociedade contemporânea ” ( que deveria independer da interpretação de \textbf{cada} grupo social a respeito da sociedade na qual vive ).
No entanto, apesar de envolver tais conceitos, espera -se que alunos concluintes de cursos superiores de graduação sejam capazes de discutir funções de instituições sociais, especialmente porque o enunciado trazia fato recente e de grande repercussão como apoio. superiores de graduação sejam capazes de discutir funções de instituições sociais, especialmente porque o enunciado trazia fato recente e de grande repercussão como apoio.
O padrão de resposta elencou as funções sociais que poderiam ser abordadas nas respostas dos estudantes, como também argumentos que poderiam ser usados ao longo das mesmas.
A tarefa dos estudantes era a construção de um texto expositivo que articulasse a importância de museus, suas funções sociais e as perspectivas de memória e de futuro, o que se mostrou de dificuldade de média a difícil.
Apesar de o enunciado ser claro, muitos estudantes, provavelmente influenciados pelos textos, limitaram -se a discutir aspectos relacionados ao Museu Nacional e, em particular, ao incêndio que o destruiu.
Foram poucos os que interpretaram corretamente o escopo do comando, que solicitava a discussão das funções sociais de museus de uma forma geral, no sentido de argumentar a importância deste tipo de instituição.
As quatro funções sociais elencadas no padrão de resposta foram encontradas na quase totalidade das redações dos estudantes de forma articulada com a importância dos museus.
Entretanto, quase nunca os textos abordavam mais do que uma função social e vinculavam, de forma clara, a importância desse tipo de instituição para a sociedade contemporânea.
Além disso, a solicitação de contextualização sob o foco museu/memória/perspectiva de futuro parece não ter sido bem compreendida e, em consequência, \textbf{esteve} pouco presente nas respostas.
Também se observou confusão entre os conceitos de história e memória e uma visão estreita do significado de sociedade contemporânea, às vezes, restrita à cidade do respondente, outras, ao Rio de Janeiro, ou ainda a um tipo de sociedade idealizada.
A relevância que os textos de apoio deram ao caso particular do incêndio do Museu Nacional do Rio de Janeiro e aos prejuízos causados ao seu acervo, induziu a maioria dos estudantes a elaborar respostas tratando das funções a partir da perda e, em particular, à perda do tipo de acervo daquele museu e do tipo de pesquisa que lá se desenvolve, divergindo do comando da questão, que solicitava texto abordando funções sociais de museus em geral.
Também foram encontradas muitas respostas que se limitavam a discutir o que o estudante considerava como possíveis causas do incêndio, a falta de manutenção, questões políticas e de gestão, sem responder ao que foi solicitado.
Foi também recorrente a transcrição de frases dos textos do enunciado nas respostas, diretamente ou por paráfrases, revelando \textbf{certa} dificuldade de criação textual.
A análise das respostas dos estudantes, além de revelar dificuldade de abordar a temática esperada, revelou outros dois aspectos predominantes.
Um deles, como já comentado, foi a ênfase no incêndio do Museu Nacional, que orientou grande parte dos textos dos estudantes.
Outro, também evidenciando desvio do esperado, foi o foco na perda patrimonial.
Nesse grupo, houve respostas que até apresentaram uma função social, mas olhada pelo ponto de vista da perda.
Nessa linha, muitas respostas enalteciam a importância do Museu Nacional, sua relevância no cenário nacional, o impacto da perda de seu acervo, isso complementado por frágeis reflexões sobre as funções sociais.
Em consequência, a perspectiva de desdobramentos mais ricos, que tratassem de aspectos sociais relacionados a diferentes tipos de museu foi prejudicada.
Muitas respostas trataram da preservação da memória e/ou guarda do acervo ; identidade, tanto individual quanto coletiva ; conhecimento sobre antepassados ; evolução da sociedade e dos seus costumes ; valorização e disseminação do conhecimento por meio da liberação do acesso aos acervos ; bem como atividades pedagógicas e exposições de arte.
O aspecto produção de pesquisa e conhecimento, embora frequente, foi mencionado de forma muito simples, sem qualquer aprofundamento, para além do que já era mencionado no enunciado.
Frequentemente apresentava -se o museu como produtor de pesquisas para o campo da saúde, o que evidencia desconhecimento das reais possibilidades de pesquisa, mesmo no Museu Nacional.
Quanto a perspectiva de futuro, quando presente no texto, era mencionada de forma breve, com afirmações genéricas sobre a possibilidade de melhoria do futuro pelo simples conhecimento do passado, numa percepção de que o passado determina o futuro, de forma linear.
Foram raras as respostas com alguma argumentação que associasse pesquisa, memória, identidade e reflexão social como base para a construção de alternativas de futuro.
A relação entre turismo e museu apareceu considerando a visita a museus como forma de lazer e, algumas vezes, como geradora de recursos.
O desempenho pode ser considerado fraco.
Muitas respostas deixavam claro que o estudante nunca tinha visitado um museu e, de uma forma geral, observou -se desconhecimento sobre suas funções sociais.
Não se esperava que história, memória social e museologia fossem assuntos dominados pelos estudantes das diversas áreas que participaram do Enade 2018, mas esperava -se que essa temática pudesse ser abordada do ponto de vista das experiências pessoais ou de conhecimento social.
Como uma vivência própria com os museus também se revelou frágil ou inexistente, o conhecimento sobre o assunto ficou bem restrito aos textos motivadores e ao senso comum, construído, basicamente, pelo noticiário em \textbf{torno} do incêndio.
O universo dos argumentos dos estudantes se mostrou muito limitado. \textbf{Cada} dimensão foi praticamente representada por uma ideia simplista.
De forma recorrente : a preservação da memória para construção da identidade de um povo ; a compreensão do passado por meio das pesquisas como oportunidade para melhorar o futuro ; o museu como um espaço em que há um aprendizado estimulado pelo lazer.
Houve, também, uma grande concentração de respostas associando o museu a um local de guarda de acervo e de história.
As respostas à Questão 2 evidenciaram dificuldades conceituais.
Além dos conceitos de história, memória e museu se misturarem, parecendo não se distinguirem, ao mencionarem a pesquisa, também se observou falta de distinção entre pesquisas colegiais, com objetivos de aprendizado, e as realizadas por pesquisadores, com objetivos de avançar cientificamente.
A falta de compreensão do que se pedia no comando da questão e a dificuldade de expressão linguística foram dois aspectos marcantes na avaliação dos corretores.
Foi frequente os textos apresentarem ideias soltas ou de forma confusa, sem constituir unidade.
O vocabulário usado se revelou bem restrito, e fortemente apoiado no coloquial, na linguagem oral.
Predominou o senso comum em lugar de uma reflexão mais formal envolvendo os conceitos envolvidos. 6.3.1.5 Análise de Língua Portuguesa das Questões Discursivas do Componente de Formação Geral Os dados de Tecnologia em Design de Interiores, obtidos a partir das respostas às questões discursivas do Componente de Formação Geral, no que tange à Língua Portuguesa, encontram -se na Tabela 6.16 e no Gráfico 6.9.
Nesse aspecto, os alunos de todo o Brasil obtiveram Média 53,2.
A maior Média, em relação à Língua Portuguesa, 54,4, foi obtida nas regiões Sul e Sudeste, e a menor, na região Norte ( 45,7 ).
Quanto à variabilidade das notas, o Desvio padrão de todo o Brasil foi 24,2.
O menor Desvio padrão foi obtido na região Sudeste ( 22,7 ) e o maior Desvio padrão foi obtido na região Centro-Oeste ( 26,9 ).
A Mediana das notas de Língua Portuguesa foi 60,0 para o Brasil, como um todo, a mesma obtida nas regiões Nordeste e Centro-Oeste.
A Mediana da região Norte foi 55,0, a menor entre as obtidas pelas Grandes Regiões, enquanto as das regiões Sudeste e Sul foram 62,5, as maiores.
A nota Máxima para todo o Brasil foi 92,5, com, pelo menos, um aluno tirando essa nota na região Nordeste.
Na região Norte, a nota Máxima foi 82,5 ( menor nota máxima entre as Grandes Regiões ) e, nas regiões Sudeste, Sul e Centro-Oeste, 90,0.
Já a nota Mínima foi zero em todas as regiões do país.
No Gráfico 6.9, mostra -se a distribuição das notas de Língua Portuguesa do Componente de Formação Geral.
Observa -se que a maior frequência ( 27,2 \% ) corresponde à dos alunos que obtiveram nota no intervalo ( 60 ; 70 ].
Destacam -se, também, os alunos que deixaram ambas as questões em branco, representando 10 \% do total. 6.3.1.6 Comentários sobre a correção das respostas de Formação Geral com respeito à Língua Portuguesa As duas questões discursivas de Formação Geral ensejaram duas configurações textuais distintas, tendo em vista os seus encaminhamentos : a Questão 1 teve um caráter mais opinativo e a Questão 2, mais expositivo, conforme se explicita a seguir.
Ambas as questões permitiram a observação do desempenho linguístico dos participantes em situação formal, competência exigida no processo de formação universitária.
As habilidades de produção textual envolvidas têm, igualmente, como condição básica, o domínio do padrão formal culto da Língua Portuguesa, no que diz respeito aos aspectos relativos às questões de caráter ortográfico, textual, morfossintático e vocabular.
O encaminhamento das questões evidencia a expectativa de que o participante utilizasse seus conhecimentos sobre os dois temas e estruturasse seus textos de acordo com as características do registro formal adequado à situação comunicativa.
Tal configuração determina exigências quanto : à adequação da seleção vocabular, ao desenvolvimento do conteúdo, à estruturação sintática dos períodos, à organização lógica das ideias, à utilização de procedimentos de encadeamento textual e de referenciação, à obediência às exigências morfossintáticas próprias da modalidade escrita da norma-padrão, ao respeito às regras ortográficas e às regras de acentuação gráfica.
Esse encaminhamento exigiu do graduando o domínio do modo de organização textual de base dissertativo-argumentativa, essencial ao processo de formação universitária.
As habilidades de produção textual envolvidas têm como condição básica o domínio do padrão formal culto da Língua Portuguesa, no que diz respeito aos aspectos relativos às questões de caráter ortográfico, textual, morfossintático e vocabular.
O enunciado da Questão 1 tomava como base três fragmentos de textos jornalísticos publicados em sites da internet, que abordavam o aumento do assassinato de minorias e de defensores de direitos humanos no Brasil, entre eles a vereadora Marielle Franco.
O comando solicitava que o participante desenvolvesse um texto dissertativo e argumentativo com as temáticas Fake News, Democracia e Direitos Humanos.
Embora não tenha havido referência explícita à tipologia textual esperada para o desenvolvimento da resposta, essa proposta encaminhou o participante para a elaboração de um texto opinativo, com o objetivo de discutir o problema e manifestar um posicionamento crítico.
O enunciado da Questão 2 abordava o incêndio do Museu Nacional do Rio de Janeiro, ocorrido em 2018, e seus impactos na ciência brasileira e internacional, tomando como base dois fragmentos de textos jornalísticos publicados em sites da internet.
O comando solicitava que o participante desenvolvesse um texto abordando três aspectos da função social dos museus na sociedade contemporânea, encaminhando -o para a elaboração de um texto expositivo que confirmasse a posição defendida nos fragmentos motivadores a respeito da importância dessas instituições.
Para atender a essas exigências, o padrão de resposta utilizado na avaliação das questões discursivas de Formação Geral considerou aspectos relevantes ao bom desempenho linguístico como competências distintas, de modo a permitir um mapeamento detalhado do domínio dos recursos disponíveis na Língua Portuguesa para a comunicação escrita formal : ortográficos, textuais, morfossintáticos e vocabulares ( ver detalhes no Anexo VIII – Padrão de Respostas ).
Aspectos Ortográficos Para avaliar a competência relativa ao domínio das convenções ortográficas da norma-padrão da Língua Portuguesa, observou -se o respeito às regras de acentuação gráfica e da grafia padrão das palavras ( com ausência de abreviaturas próprias da linguagem da internet ), de acordo com as convenções estabelecidas pela legislação em vigor e consubstanciadas no Vocabulário Ortográfico da Língua Portuguesa, editado pela Academia Brasileira de Letras.
Aspectos Textuais Para avaliar a competência relativa ao domínio dos procedimentos de estruturação textual, procurou -se observar a utilização de mecanismos linguístico-discursivos que estabeleçam o encadeamento lógico entre as partes do texto, de forma a garantir a progressão e a coerência textuais.
Esta competência destaca -se por sua relevância, já que são avaliadas a clareza e a unidade textual.
Aspectos Morfossintáticos Para avaliar o domínio dos diferentes aspectos morfossintáticos próprios da modalidade escrita formal da norma-padrão da Língua Portuguesa, observou -se se o texto produzido atendeu às seguintes exigências : a concordância nominal, a concordância verbal, a regência nominal, a regência verbal, a flexão nominal, a flexão verbal, a correlação entre os tempos verbais, a colocação pronominal e a utilização de sinais de pontuação que contribuam para a organização lógica da frase e do texto.
Aspectos Vocabulares Para avaliar a adequação da seleção vocabular à modalidade escrita formal da Língua Portuguesa, exigida pela situação comunicativa, observou -se o atendimento às seguintes exigências : precisão na seleção/utilização do vocabulário relacionado à temática solicitada pela questão ; ausência de marcas da oralidade, como termos de sentido muito genérico ( “ coisa ”, “ negócio ”, “ você ” ) e termos de registros mais informais ( como gírias, jargões, frases feitas, ditados populares, termos regionais ). ( “ coisa ”, “ negócio ”, “ você ” ) e termos de registros mais informais ( como gírias, jargões, frases feitas, ditados populares, termos regionais ).
Assim, espera -se que o participante respeite a adequação vocabular não usando gírias ou expressões coloquiais, evite repetição desnecessária de palavras e utilize um vocabulário mais formal, como solicitado por um texto dissertativo.
Para efeito de pontuação, a avaliação do desempenho linguístico considerou três grandes grupos de competências, agrupando aspectos explicitados, com os seguintes pesos relativos : aspectos ortográficos ( 20 \% ) ; aspectos textuais ( 40 \% ), aspectos morfossintáticos e vocabulares ( 40 \% ).
Os critérios de pontuação utilizados no padrão de resposta procuraram atender à diversidade de desempenho dos graduandos.
Nos textos analisados, observaram -se, assim, vários níveis de atendimento às exigências do padrão formal próprio da modalidade escrita da Língua Portuguesa, configurando um"continuum".
Por um lado, textos que traduzem com exatidão a expectativa de uso formal da língua, em todos os níveis analisados, marcados por clareza e unidade ; e, no outro extremo, textos quase incompreensíveis devido ao grande número de desvios de todos os tipos, principalmente os de caráter textual.
As duas questões discursivas de Formação Geral analisadas revelaram desempenhos distintos dos participantes.
A Questão 1 propiciou a oportunidade de maior desenvolvimento do tema solicitado, gerando consequências na estruturação textual, já que os textos foram mais longos, mais elaborados e fluentes.
Em virtude da polêmica sobre a existência das fake news e sobre a visão pública da personagem da vereadora Marielle Franco, observaram -se respostas que polemizaram o tema, favorecendo a construção argumentativa do texto.
A Questão 2 teve um comportamento distinto, devido ao caráter mais expositivo do enunciado e ao reduzido conhecimento dos participantes sobre o tema.
Assim, os textos foram mais curtos, frequentemente com repetição de ideias dos textos motivadores.
Quanto aos aspectos linguísticos analisados durante esta avaliação, observaram -se os seguintes resultados : Aspectos ortográficos : O desempenho dos participantes revelou uma diferença muito grande nos dois aspectos analisados nesta competência : baixo índice de desvios da grafia padrão e grande índice de desvios de acentuação.
Em vários casos, ocorre ausência completa de acentuação gráfica em todas as palavras do texto.
Os resultados revelaram que a tendência dominante entre os universitários brasileiros é a eliminação da acentuação gráfica, talvez motivada pelos hábitos relacionados às redes sociais e pela ausência de esclarecimento dos meios de comunicação, das autoridades e das escolas sobre as decisões do Acordo Ortográfico da Língua Portuguesa de 1990.
Os casos mais sistemáticos de eliminação do acento indicador da sílaba tônica foram : - palavras proparoxítonas ( por exemplo “ generos ”, “ arbitro ”, “ politico ”, “ publicas ”, “ numero ” ) ; - palavras paroxítonas terminadas em ditongo crescente ( por exemplo “ area ”, “ noticiarios ”, “ individuo ”, “ historia ”, “ varios ”, “ noticia ”, “ negligencia ”, “ ocorrencia ” ) ; - palavras oxítonas ( por exemplo “ ninguem ”, “ esta ”, “ ate ”, “ tambem ” ).
Por outro lado, destaca -se o uso indevido do acento gráfico em palavras como “ melâncolia ” e “ intervênção ”, por exemplo.
Quanto ao domínio das convenções relativas à grafia das palavras, observaram -se desvios esporádicos como : hipercorreção pela escolha de “ e ” no lugar de “ i ”, por influência de hábitos da oralidade ( “ descriminação ” no lugar de “ discriminação ”, “ entervenção ” no lugar de “ intervenção ” ) ; desvios de grafia relacionados à variação diastrática, como por exemplo “ perca ”, “ subjulgadas ”, “ precoseito ”, “ soubre ”, “ vecendo ”, “ indesencia ”, “ apolojia ”, “ fulga ”, “ dereitos ” ;...
Observaram -se, também, muitos casos de inadequação no uso da maiúscula : uso indevido para destacar determinadas palavras-chave do texto, como “ Brasileiros ”, “ Homens ” ; ausência de maiúscula para grafar nomes próprios, como"marielle".
Expressões como"Direitos Humanos"foram grafadas diversificadamente, com maiúsculas e com minúsculas.
Destaque -se, também, a maioria dos participantes que grafam os textos inteiramente em caixa alta, o que impede a identificação de desvios relacionados ao uso de maiúsculas.
Vale observar, também, que, ao contrário do que se esperava, não apareceram abreviaturas próprias do “ internetês ”, relacionadas ao uso de redes sociais e e-mails, como por exemplo"vc".
Aspectos textuais : Para alguns dos graduandos essa competência se revelou como a mais problemática, tendo em vista os inúmeros problemas observados, desvios acumulados durante toda a formação escolar.
São eles : sequência justaposta de ideias sem encaixamentos sintáticos, provocando uma redução de estruturas subordinadas, ao lado do aumento na frequência de estruturas coordenadas e absolutas ; frases fragmentadas que comprometem a estrutura lógico-gramatical ; frases formadas apenas por oração subordinada, sem oração principal ; redução no uso de conectores para expressar relações lógicas essenciais à construção do texto como consequência da mudança de estruturação frasal ; emprego inadequado do pronome relativo ( com omissão da preposição ou a utilização do pronome relativo inadequado, como “ onde ” ), refletindo hábitos da oralidade ; ausência de recursos de referenciação, como substituição de termos por sinônimos, hiperônimos, nominalizações, expressões metafóricas.
Esses problemas revelam dificuldades em relação à estrutura formal do texto produzido, o que é preocupante ao se levar em conta que são graduandos em fase final de formação.
Em alguns textos, falta um mínimo de textualidade e de domínio do registro padrão da língua.
Na verdade, observaram -se relações linguísticas quase agramaticais, como as estabelecidas pela sequência de gerúndios sem o apoio de um ponto de partida para a organização das informações gramaticais e semânticas.
Quanto à utilização dos sinais de pontuação, observou -se uma ausência dos mesmos na maioria dos textos analisados.
Foi muito frequente a ocorrência de parágrafos sem marca interna de pontuação para separar os períodos.
São os seguintes os tipos de problemas mais encontrados :  vírgula : utilização de vírgula para separar o sujeito e o predicado ; uso de vírgula no lugar do ponto para separar ideias que constituem períodos distintos ; ausência de vírgula para separar oração adjetiva explicativa ;  ponto e vírgula : utilização do ponto e vírgula no lugar de vírgula ;  ponto final : ausência de ponto final para separar períodos.
Aspectos morfossintáticos e vocabulares : Os resultados são transparentes em relação aos aspectos mais problemáticos no desempenho dos participantes nestas duas competências.
Em relação aos aspectos morfossintáticos, seguem algumas observações sobre os desvios mais frequentes.
Quanto à regência, o desvio mais frequente foi a falta do sinal indicativo da crase – isso revela que o usuário não tem consciência de que, sob a forma do termo “ a ”, existe a presença de uma contração entre a preposição “ a ” ( exigida pela regência do termo anterior ) e o artigo definido “ a ”.
Outro problema relacionado à regência verbal e à nominal, encontrado frequentemente nas respostas, foi a ausência de preposição antes de pronome relativo, processo generalizado lógico-gramatical ; frases fragmentadas que comprometem a estrutura lógico-gramatical ; frases formadas apenas por oração subordinada, sem oração principal ; redução no uso de conectores para expressar relações lógicas essenciais à construção do texto como consequência da mudança de estruturação frasal ; emprego inadequado do pronome relativo ( com omissão da preposição ou a utilização do pronome relativo inadequado, como “ onde ” ), refletindo hábitos da oralidade ; ausência de recursos de referenciação, como substituição de termos por sinônimos, hiperônimos, nominalizações, expressões metafóricas.
Esses problemas revelam dificuldades em relação à estrutura formal do texto produzido, o que é preocupante ao se levar em conta que são graduandos em fase final de formação.
Em alguns textos, falta um mínimo de textualidade e de domínio do registro padrão da língua.
Na verdade, observaram -se relações linguísticas quase agramaticais, como as estabelecidas pela sequência de gerúndios sem o apoio de um ponto de partida para a organização das informações gramaticais e semânticas.
Quanto à utilização dos sinais de pontuação, observou -se uma ausência dos mesmos na maioria dos textos analisados.
Foi muito frequente a ocorrência de parágrafos sem marca interna de pontuação para separar os períodos.
São os seguintes os tipos de problemas mais encontrados :  vírgula : utilização de vírgula para separar o sujeito e o predicado ; uso de vírgula no lugar do ponto para separar ideias que constituem períodos distintos ; ausência de vírgula para separar oração adjetiva explicativa ;  ponto e vírgula : utilização do ponto e vírgula no lugar de vírgula ;  ponto final : ausência de ponto final para separar períodos.
Aspectos morfossintáticos e vocabulares : Os resultados são transparentes em relação aos aspectos mais problemáticos no desempenho dos participantes nestas duas competências.
Em relação aos aspectos morfossintáticos, seguem algumas observações sobre os desvios mais frequentes.
Quanto à regência, o desvio mais frequente foi a falta do sinal indicativo da crase – isso revela que o usuário não tem consciência de que, sob a forma do termo “ a ”, existe a presença de uma contração entre a preposição “ a ” ( exigida pela regência do termo anterior ) e o artigo definido “ a ”.
Outro problema relacionado à regência verbal e à nominal, encontrado frequentemente nas respostas, foi a ausência de preposição antes de pronome relativo, processo generalizado na modalidade oral da língua, em situações de registro informal.
Apesar da possibilidade de que essa alteração de regência se generalize no padrão escrito da Língua Portuguesa, como já \textbf{está} ocorrendo até em textos jornalísticos, o não emprego da preposição foi considerado como desvio neste processo de avaliação.
Outro desvio de regência significativo foi a utilização inadequada de uma preposição ou sua ausência após o verbo ou o nome ( substantivo ou adjetivo ).
A concordância verbal e a concordância nominal apresentaram alguns desvios muito frequentes.
Quanto à concordância de número, observou -se, como apontam as pesquisas nessa área, ausência de marca ( com sujeito anteposto ou posposto ) ou uso indevido ( uso inadequado da marca de plural comandado pelo núcleo plural da locução adjetiva, apesar de o substantivo que funciona como núcleo do sintagma nominal \textbf{estar} no singular ).
Uma ocorrência que se destacou foi a ausência de acento circunflexo na forma plural do presente do indicativo do verbo “ ter ”, que foi considerada como um desvio na concordância verbal e não na acentuação gráfica.
Quanto à concordância de gênero, vários casos foram observados, normalmente no âmbito de sintagmas nominais longos, em que o adjetivo \textbf{está} afastado do substantivo.
Quanto à questão da colocação pronominal, foram poucos os desvios observados.
Concluiu -se que, no registro escrito formal, a maioria dos participantes já incorporou regras como a não introdução da frase por um pronome oblíquo e a próclise na presença de um termo atrator.
Várias marcas de oralidade foram identificadas, embora não com alta frequência : uso do pronome relativo “ onde ” como relativo universal, falta de artigo definido antes de substantivo isolado ou antes de termos em enumerações.
Quanto aos aspectos vocabulares, alguns tipos de inadequação foram observados : expressões da oralidade ; seleção vocabular incompatível com o contexto, gerando falta de inteligibilidade ; falta de domínio de vocabulário mais abstrato e de maior complexidade, essencial ao desenvolvimento do texto de base dissertativa.
O principal aspecto observado foi a excessiva repetição de \textbf{certas} palavras, como o termo"pessoa", por exemplo, revelando limitação de repertório vocabular.
Em função do tema solicitado na Questão 1, diferentes grafias da expressão fake news não foram consideradas.
Para sintetizar, os problemas que mais se destacaram nos textos analisados foram : a ) truncamentos sintáticos, caracterizados por estruturas frasais incompreensíveis devido à ruptura da complexidade sintática própria da modalidade escrita ; b ) redução de estruturas subordinadas, compensada pelo aumento na frequência de estruturas coordenadas e absolutas, por um lado, ou pela elaboração de estruturas textuais prejudicadas pelo excesso de ideias sem a devida conexão subordinativa ; c ) redução no uso de conectores para expressar relações lógicas essenciais à construção do texto, substituídos pela exigência de inferência por parte do interlocutor para suprir a sua ausência ; d ) redução no uso do subjuntivo, ao lado da ampliação do uso do indicativo, combinado a estruturas frasais coordenadas ou absolutas ; d ) empobrecimento do processo de referenciação, com a repetição de palavras, sem processos mais sofisticados de substituição ; e ) simplificação da marcação da categoria tempo na morfologia verbal ; f ) desrespeito às regras de concordância verbal e nominal próprias do padrão formal ; g ) ausência do acento grave indicador da crase, revelando falta de conhecimento dos critérios de regência verbal e nominal ; h ) falta de domínio de vocabulário mais abstrato e de maior complexidade, essencial ao desenvolvimento do processo dissertativo ; i ) redução do emprego da acentuação gráfica.
Como resultado da correção, merece atenção a constatação de que a modalidade escrita tem apresentado a tendência a uma evidente simplificação, aproximando -se das características da modalidade oral da Língua Portuguesa.
No caso dos textos de base dissertativa ( expositiva ) e de base argumentativa, inscritos em um registro formal, a distância entre as duas modalidades é ainda maior, o que provoca desvios recorrentes em todos os quatro aspectos analisados : ortográficos, textuais, morfossintáticos e vocabulares.
Destaca -se, nessa avaliação, o comprometimento textual, por meio de estruturas fragmentadas e/ou truncadas, rompendo a complexidade sintática esperada no padrão formal. 6.3.2 Componente de Conhecimento Específico Na parte da prova relativa às questões discursivas do Componente de Conhecimento Específico ( Tabela 6.17 ), observa -se que a Média foi mais baixa do que para as questões discursivas do Componente de Formação Geral.
Enquanto no Componente de Formação Geral a Média para estudantes de Tecnologia em Design de Interiores de todo o Brasil foi 35,6, na parte de Conhecimento Específico, a Média foi 28,1.
A maior Média deste componente foi obtida pelos estudantes da região Sul ( 29,9 ), e a menor, pelos da região Norte ( 22,9 ).
Quanto à variabilidade das notas, o Desvio padrão de todo o Brasil foi 15,8.
O maior Desvio padrão foi encontrado na região Norte ( 17,0 ), e o menor, nas regiões Sudeste e Centro-Oeste ( 15,4 ).
A maior nota Máxima, 80,0, foi obtida na região Sudeste.
Nas demais regiões, as notas Máximas foram : 68,3, na regiões Norte, 76,7, na Nordeste, 75,0, na Sul, e 71,7 na região Centro-Oeste.
A nota Mínima ( 0,0 ) foi obtida por, pelo menos, um aluno em todas as regiões do Brasil.
A Mediana do Brasil, como um todo, foi 28,3.
A maior Mediana foi obtida na região Sul ( 31,7 ), e a menor, na região Norte ( 21,7 ).
No Gráfico 6.10, representa -se a distribuição das notas nas questões discursivas no Componente de Conhecimento Específico.
A moda dessa distribuição ocorre no intervalo [ 20 ; 30 ], com mais do que 23 \% do total de participantes.
O intervalo [ 0 ; 10 ], no qual \textbf{estão} computadas as ocorrências de respostas em branco para as três questões discursivas de Conhecimentos Específicos, teve frequência de cerca de 16 \% do total de participantes.
A análise de \textbf{cada} uma destas questões será feita a seguir. 6.3.2.1 Análise da Questão Discursiva 3 do Componente de Conhecimento Específico Na questão 3, cujos resultados aferidos encontram -se descritos na Tabela 6.18, a Média dos estudantes de todo o Brasil foi 32,9, a questão de melhor desempenho dentre as três discursivas de Conhecimentos Específicos.
A menor Média nessa questão foi obtida pelos alunos da região Centro-Oeste ( 26,2 ), enquanto a maior Média foi obtida na região Sul ( 37,1 ).
Quanto à variabilidade das notas, o Desvio padrão de todo o Brasil foi 26,1.
O maior Desvio padrão foi obtido na região Sul ( 26,8 ), enquanto o menor foi obtido na região Centro-Oeste ( 24,9 ).
A nota Máxima, 100,0 pontos, foi alcançada por, pelo menos, um aluno da região Sudeste e um da região Sul.
As demais notas máximas foram : 85,0, na região Norte ; 90,0, na região Nordeste ; e 95,0, na região Centro-Oeste.
A Mediana do Brasil, como um todo, foi 35,0, a mesma das regiões Nordeste e Sudeste.
Nas demais regiões, as medianas foram : 30,0, na região Norte ; 40,0, na região Sul ; e 25,0, na região Centro-Oeste.
A nota Mínima ( 0,0 ) também foi a mesma em todas as regiões do Brasil.
No Gráfico 6.11, mostra -se a distribuição das notas na questão discursiva 3, do Componente de Conhecimento Específico da área de Tecnologia em Design de Interiores.
A moda dessa distribuição é a dos alunos que receberam nota zero, correspondendo a 14,9 \%.
A segunda maior frequência, 12,5 \%, é no intervalo ( 30 ; 40 ], seguida pelos intervalos ( 40 ; 50 ], com frequência igual a 12,1 \%, e ( 20 ; 30 ], com frequência 11,4 \%.
No Gráfico 6.11, mostra -se a distribuição das notas na questão discursiva 3, do Componente de Conhecimento Específico da área de Tecnologia em Design de Interiores.
A moda dessa distribuição é a dos alunos que receberam nota zero, correspondendo a 14,9 \%.
A segunda maior frequência, 12,5 \%, é no intervalo ( 30 ; 40 ], seguida pelos intervalos ( 40 ; 50 ], com frequência igual a 12,1 \%, e ( 20 ; 30 ], com frequência 11,4 \%. 6.3.2.2 - Comentários sobre as respostas à Questão Discursiva 3 A Questão 3, claramente pertinente às Diretrizes Curriculares da área de Design de Interiores, trazia um texto introdutório que descrevia os arranjos físicos para ambientes de escritórios que seguem o padrão de Coworking, além de uma foto de um escritório nos moldes do desse padrão.
O Coworking é uma tipologia que abrange o compartilhamento do local de trabalho por indivíduos ou sociedades que não atuam, necessariamente, na mesma área profissional.
Esses espaços comuns são organizados em layouts específicos, e a funcionalidade e eficiência do projeto dependem do equilíbrio das diretrizes de conforto ambiental.
O comando da questão solicitava duas soluções técnicas de design para \textbf{cada} área de conforto ambiental ( térmico, lumínico e acústico ) além da explicação como essas soluções técnicas favorecem o desempenho funcional e a eficiência energética do local, possibilitando a expressão do domínio do estudante sobre a matéria.
Como o enunciado não definiu o teor das soluções—se passivas, ativas ou ambas —, assim como, se intervenções pertinentes ao envoltório arquitetônico poderiam ser explicitadas pelos estudantes, na correção todas elas foram aceitas : entendeu -se que a Lei 13.369/16 de regulamentação da profissão de Design de Interiores admite a proposição, por parte dos designers, de intervenção estruturais, a serem validadas por profissional devidamente habilitado.
A complexidade da questão \textbf{está} na necessidade de esquematização da resposta para que nenhuma das abordagens seja esquecida, pois são nove os itens a serem apresentados — quais sejam, duas soluções para \textbf{cada} um dos três tipos de conforto, acrescidas das justificativas de seu uso através de aspectos funcionais e econômicos.
A divisão da pontuação entre os nove itens determinou a maior valorização das soluções passivas de Conforto Lumínico e de Conforto Térmico, compreendendo -se a sua relevância para o desenvolvimento sustentável, especificamente na área de Interiores.
A questão pode ser considerada fácil, sabendo -se que conteúdos técnicos são, muitas vezes, mais bem assimilados pelos estudantes.
A presença de uma imagem ilustrativa no enunciado possibilitou o seu uso para inferência de respostas, como facilmente observado em alguns casos.
Igualmente verificou -se resposta ao conforto ambiental como um todo, sem discriminação do seu tipo, dando indícios de incompreensão do enunciado.
A resposta incompleta foi a mais comum, talvez pela dificuldade de organização da mesma, pois o comando era claro.
Houve ainda ocorrência, com alguma frequência, de apresentação de número superior a duas especificações corretas para apenas um dos itens de conforto, e nenhuma para outro, como solicitado.
Embora algumas dessas soluções pudessem também servir, eventualmente, para outro tipo de conforto, foi considerado necessário o apontamento a essa segunda função para que houvesse aproveitamento da pontuação.
Nem sempre a linguagem foi clara, percebendo -se o acerto da solução sem que houvesse uso de terminologia técnica pertinente.
Esse fato evidencia alguma insuficiência no uso de uma linguagem adequada.
De modo geral, houve baixa diversidade de soluções, sendo raras aquelas que fugiam ao padrão, além de um grande número de respostas incompletas e de um número reduzido de respostas que separavam, de modo organizado, as soluções apresentadas.
Quanto à explicitação de uma única opção técnica para atender a mais de um item — como, por exemplo, a indicação de"grandes janelas", que poderiam atender tanto ao conforto térmico quanto ao lumínico, — exigiu -se que a resposta deixasse evidente as duas perspectivas de uso.
Das três questões discursivas, essa foi a que obteve a maior média, que, embora baixa, evidencia maior facilidade.
No entanto, o desempenho observado pode ser considerado insatisfatório, por se tratar de conteúdo técnico básico para a atuação na área de Design de Interiores.
Respostas consideradas fracas e médias foram a grande maioria, aliadas a um pequeno número de respostas boas.
Essa distribuição de notas evidencia um ponto de fragilidade na formação dos estudantes.
Tais números mostram uma distribuição concentrada no resultado inferior, identificando uma falha relevante para a área na assimilação de conteúdo.
O número ampliado de quesitos ( nove ) a serem respondidos, somando -se às dificuldades dos estudantes na organização da resposta e na dificuldade de expressão escrita, foi o aspecto mais observado.
Não há dúvidas de que, embora com melhor desempenho na média entre as três questões discursivas, os resultados mostram -se aquém do esperado, para uma questão de \textbf{ordem} técnica que deveria ser de domínio geral.
Os resultados apontam dificuldade na expressão escrita e na assimilação de uma linguagem técnica adequada.
Pode -se apontar para um conhecimento que, embora apresentado aos estudantes, não foi totalmente assimilado por falta de materialização do ato projetual.
No entanto, em termos teóricos, não há dúvida de que há um déficit na formação dos estudantes. 6.3.2.3 - Análise da Questão Discursiva 4 do Componente de Conhecimento Específico Na Tabela 6.19, constam as informações relativas à questão 4 do conjunto de questões do Componente de Conhecimento Específico.
O desempenho dos estudantes nessa questão foi inferior ao da questão 3 e superior ao da questão 5.
A Média geral do Brasil foi 26,6, sendo a maior Média registrada na região Sudeste ( 28,2 ), e a menor, na região Norte ( 21,1 ).
A nota Máxima ( 100,0 ) foi atingida por, pelo menos, um concluinte da região Sudeste.
Nas regiões Nordeste e Sul, a nota Máxima foi 90,0, e, nas regiões Norte e Centro-Oeste, 80,0.
A Mediana em todo o Brasil foi 30,0, o mesmo valor foi obtido nas regiões Sudeste e Sul.
Nas regiões Norte e Centro-Oeste, a Mediana foi igual a 20,0, e, na Nordeste, 25,0.
A nota Mínima ( 0,0 ) também foi a mesma em todas as regiões do Brasil.
No Gráfico 6.12, \textbf{está} representada a distribuição de notas da questão discursiva 4, do Componente de Conhecimento Específico.
Essa distribuição tem moda no intervalo ( 30 ; 40 ], com frequência 18,1 \%, seguido de perto pelas classes de estudantes que deixaram a questão em branco, com frequência de 17,3 \%, e também pela classe dos que \textbf{estavam} no intervalo ( 20 ; 30 ], com frequência 17,0 \%.
A soma das frequências destes três intervalos supera 50 \% dos respondentes. 6.3.2.4 Comentários sobre as respostas à Questão Discursiva 4 A questão discursiva 4 apresentava um texto introdutório e fotos de duas salas de jantar decoradas em períodos distintos, ambas com inspiração no Neoclassicismo.
O texto introdutório se referia ao surgimento do Neoclassicismo em meados do século XVIII e o contexto histórico da relação do ser humano com os novos conceitos da era industrial, da arte e da arquitetura, que nesse momento buscava um novo estilo.
O texto abordava também a difusão do Neoclassicismo pela Europa com a publicação de livros e manuais.
O comando da questão solicitava que o aluno citasse quatro características gerais do Neoclassicismo, além de elaborar uma análise comparativa dos ambientes, destacando os princípios de composição espacial.
A Questão 4 é pertinente às diretrizes curriculares da área de Tecnologia do Design de Interiores, e com isso foi possível ao estudante evidenciar seus conhecimentos sobre a Teoria e História da Arte, História do Design e do Mobiliário.
Era esperado do aluno, além da explicitação de quatro características gerais do Neoclassicismo ( item ‘ a ’ ), que também aplicasse os seus conhecimentos na elaboração de uma análise comparativa de duas imagens disponibilizadas ( item ‘ b ’ ), em que fossem destacados os princípios de composição espacial, os paralelismos e as diferenças.
As exigências do enunciado são claras, fazendo -se apenas uma sinalização quanto ao uso da expressão"características gerais", permitindo que o olhar se voltasse tanto a aspectos gerais do movimento quanto ao vocabulário formal específico do mobiliário e dos espaços interiores neoclássicos.
Em ambos os casos, e a partir das imagens apresentadas, o foco nos espaços interiores pode contribuir para o entendimento de que se tratasse de conhecimento aplicado a esses ambientes.
O espaço da resposta foi adequado e a questão pode ser considerada a mais difícil em relação às demais questões discursivas, tratando -se de conteúdo mais teórico que o estudante tende a ver como mais distante de sua realidade.
Na grande maioria das respostas dos alunos, havia uma mera descrição dos ambientes das imagens, sem que se empreendesse, de fato, uma análise comparativa, observando -se, de modo geral, muita dificuldade na elaboração de uma reflexão analítica.
Foi recorrente o apontamento da diferença de metragem quadrada entre os dois ambientes como item comparativo, o que, embora verdadeiro, não atendia ao que se pretendia analisar.
Um fato comum entre as respostas dos alunos foi a indicação de algumas características do estilo no item ‘ a ’ e, em seguida, em relação ao item ‘ b ’, a descrição de apenas uma das imagens, sem menção à outra.
Nesse caso, os aspectos apresentados ao item ‘ a ’ eram então repetidos.
O inverso também ocorria : o item ‘ a ’ não era respondido, mas, quando o estudante elaborava sua resposta ao item ‘ b ’, enumerava características do Neoclassicismo na mera descrição da imagem, sem a análise comparativa solicitada, seja no que diz respeito ao paralelismo entre as mesmas, seja em relação às diferenças estilísticas observadas.
Outra resposta frequente era a que apenas mencionava alguns itens, sem qualquer critério de organização no sentido de apontar a que parte se referia, pois, se em um primeiro momento apenas mencionava uma característica do estilo, no subsequente já havia clara referência a uma das imagens.
Também merece destaque a resposta — embora rara, — que mencionava o paralelismo na composição espacial dos dois ambientes, assinalando o princípio de simetria da disposição do mobiliário, sem que se referisse a quaisquer características do Neoclassicismo presentes, ou sem que houvesse resposta ao item ‘ a ’.
Em acréscimo à linguagem muitas vezes pouco adequada à descrição dos “ estilemas ” referentes ao estilo em pauta e a recorrentes erros ortográficos, observou -se muita dificuldade na expressão escrita e o uso de expressões como"aconchegante","sofisticado"e"belo"para a análise comparativa dos ambientes.
A impressão que se tem é de que a resposta é considerada longa pelos estudantes e que nem todos têm disponibilidade para respondê -la com afinco.
Em alguns casos, a resposta a um deles é pertinente e completa, enquanto a outra fica incompleta ou é respondida de modo sintético, monossilábico, sem que se possa inferir a compreensão do que se deseja sobre o estilo.
Entre os equívocos cometidos nas soluções apresentadas \textbf{estão} a incompletude das questões, respondendo -se a apenas um dos itens, muitas vezes ainda de modo insuficiente, assim como a confusão entre"características gerais"do Neoclassicismo e a indicação de aspectos do contexto histórico, político e cultural em que esse estilo surgiu.
Não se observou muita diversidade nas respostas dos alunos, havendo inclusive uma recorrência de respostas que mencionavam o \textbf{esquema} cromático predominante no estilo, o uso da madeira e o trabalho detalhado de paredes e tetos.
De modo geral as respostas foram fracas, havendo número considerável de respostas em branco.
O percentual alto de respostas fracas é um indicativo do grau de dificuldade dos estudantes.
Durante a correção da questão foi possível observar a dificuldade de expressão e a não assimilação de uma linguagem adequada aos estilemas do movimento em pauta.
Isso prejudicou o desempenho dos estudantes que fizeram uso, em grande parte, de termos inadequados para descrever uma dada característica.
Percebe -se ainda uma generalizada dificuldade na elaboração de uma análise comparativa, embora o enunciado seja claro.
A listagem de tópicos sumários e a ausência de uma explicitação das diferenças e similitudes observadas nos espaços apresentados parecem evidenciar repertório reduzido de expressão textual.
A média final da questão foi muito baixa, inferior à da questão discursiva anterior, o que pode ser explicado pelo conteúdo mais teórico do que prático nela presente.
Deve -se ratificar sempre a dificuldade na expressão escrita e no uso da terminologia adequada, aspectos presentes na grande maioria das provas. 6.3.2.5 Análise da Questão Discursiva 5 do Componente de Conhecimento Específico Na Tabela 6.20, constam as informações relativas à questão 5 do conjunto do Componente de Conhecimento Específico.
O desempenho dos estudantes de todo o Brasil nesta questão foi o mais baixo dentre as três questões discursivas desse componente.
A nota Média dos estudantes de todo o Brasil foi 24,8.
A maior Média foi registrada na região Nordeste ( 27,5 ), enquanto a menor Média foi registrada na região Norte ( 19,2 ).
Quanto à variabilidade das notas, o Desvio padrão dos alunos do Brasil, como um todo, foi 16,2.
Enquanto o maior desvio foi encontrado na região Nordeste ( 17,7 ), o menor foi encontrado na região Sul ( 15,0 ).
A Mediana para o Brasil foi 25,0, a mesma das regiões Nordeste e Sul.
Nas demais regiões, a Mediana foi 20,0.
Para o conjunto de alunos de Tecnologia em Design de Interiores do Brasil, a nota Máxima foi 80,0, obtida nas regiões Nordeste, Sudeste e Centro-Oeste, sendo 75,0 a nota máxima nas regiões Norte e Sul.
A nota Mínima do Brasil foi zero, o mesmo valor encontrado em todas as regiões.
No Gráfico 6.13, é representada a distribuição das notas da questão discursiva 5 do Componente de Conhecimento Específico.
Essa distribuição tem como modal o intervalo ( 20 ; 30 ], seguido de perto pelo intervalo ( 10 ; 20 ].
A soma da frequência desses dois intervalos corresponde a quase metade dos participantes. 6.3.2.6 Comentários sobre as respostas à Questão Discursiva 5 A Questão 5 apresentava duas imagens : a primeira era um diagrama de planejamento, e a segunda era a planta final de um projeto para showroom e para os escritórios de uma companhia de moda.
O texto introdutório da questão abordava a temática do processo da organização espacial de um projeto de interiores, e que esse processo se inicia na primeira entrevista com o cliente, estendendo -se até a sua finalização.
O comando da questão solicitava a descrição das etapas a serem seguidas no desenvolvimento de um projeto de interiores, desde a entrevista com o cliente até a execução do projeto.
A questão discursiva tem total aderência às diretrizes curriculares da área, focando na prática de projeto, tema sobre o qual todos os concluintes deveriam expressar total domínio.
Apesar do fraco desempenho geral, a questão não é difícil.
Destaque -se que pode ter havido dificuldade na interpretação do enunciado.
O enunciado cita a"entrevista"como o início dos trabalhos, da mesma forma que usa o termo"execução do projeto", entendendo -se com ele a sua materialização — o que deixa margem à interpretação de que essas fases \textbf{estariam} incluídas no tema solicitado :"etapas a serem seguidas no desenvolvimento de um projeto de interiores".
Em acréscimo, o título da segunda imagem apresentada —"planta final", contribui para a dificuldade de compreensão do enunciado, pois a imagem assim referida não condiz com um desenho de finalização. serem seguidas no desenvolvimento de um projeto de interiores".
Em acréscimo, o título da segunda imagem apresentada —"planta final", contribui para a dificuldade de compreensão do enunciado, pois a imagem assim referida não condiz com um desenho de finalização.
O fato fica ainda mais complexo quando se constata a ausência de uso de uma linguagem padronizada na área de Interiores no que diz respeito à metodologia de projeto, que serve de terminologia de disciplinas afins.
Isso dificulta a compreensão e gera equívocos de interpretação, justificando, dessa forma, a baixa média na questão.
A maioria dos alunos focou grande parte de suas respostas nos itens já citados no enunciado, quais sejam, a entrevista com o cliente e a execução, resumindo todas as etapas restantes e intermediárias, a croquis, layout ou projeto.
Em muitos casos, a linguagem utilizada pelos estudantes não foi adequada.
Foram observados recorrentes erros de grafia em palavras comumente usadas na área do Design de Interiores como croquis ( croks ), briefing ( breaffing, brienfing ), execução ( “ execulção"), anteprojeto ( “ anti projeto ” ).
Em acréscimo, verificou -se a aplicação de alguns termos fora do seu contexto correto, como"levantamento"da perspectiva, referindo -se ao movimento de levantar as paredes no software 3D.
Na verdade, o termo se aplicaria ao levantamento do espaço físico, arquitetônico ou não, referindo -se ao momento de coleta de dados para a elaboração do projeto e não à etapa de representação gráfica do projeto.
Os alunos confundiram as fases de um projeto e os tipos de desenhos necessários a \textbf{cada} uma delas — como planta baixa, corte, vista, maquete eletrônica — quando o solicitado eram etapas de projeto.
É possível que a prática pedagógica — que cria etapas de entregas de desenhos nas disciplinas de projeto — tenha feito com que o aluno entendesse que deveria haver uma \textbf{ordem} na entrega de pranchas, sem, contudo, relacioná -las à evolução do nível de solução do projeto.
O comando da questão solicitava a descrição de \textbf{cada} uma das etapas a serem seguidas no desenvolvimento de um projeto de interiores.
Contudo, na maioria das vezes, os alunos somente citaram as etapas, sem descrevê -las.
Os alunos na maioria das vezes confundiam os termos"programa de necessidades","briefing"e"conceito de projeto", evidenciando fortes discrepâncias quanto ao entendimento desses termos.
Briefing e entrevista muitas vezes foram empregados como sinônimos, o que sob nenhum ponto de vista pode ser dado como acertado.
A fase de elaboração do briefing — com suas diversas subdivisões, que englobam desde a coleta de dados, pesquisas, análises e diagnoses, estudo das normas vigentes e demais determinantes físicos, culturais, além das particularidades da função e dos usuários, que abrem caminho para o conceito de projeto — deve ser formulada como a base que sustenta o projeto.
Além da dificuldade no emprego de uma terminologia adequada, os resultados foram prejudicados pela massiva ausência de uma fase significativa na metodologia de projeto ( a do briefing ) que deveria ser explicitado em sua integridade, com vistas à definição de um conceito de projeto.
Outra confusão relevante foi entre o projeto executivo — que envolve ampliações e \textbf{detalhamentos} construtivos, com as devidas informações e especificações necessárias — e a própria execução do projeto, ou seja, sua materialização através da sua implantação, muitas vezes onde a obra se faz necessária.
Parte significativa dos alunos citou temas relativos a contratos, propostas e honorários, o que demonstra uma formação próxima da prática.
A dificuldade de expressão é evidente, com textos sem coerência e coesão, erros ortográficos e de concordância.
O resultado insatisfatório nas notas da questão discursiva 5 é incompatível com o nível de dificuldade da questão, que pode ser considerada de nível fácil.
A média dessa questão foi a mais baixa dentre as três questões discursivas.
O fato é surpreendente, pois o conteúdo não apenas é pertinente, mas básico ao ato projetual do designer de interiores.
Algumas respostas sugerem uma formação que extrapola os limites de atuação dos designers de interiores, alcançando os limites com a arquitetura, quando, por exemplo, são citados, como etapas do projeto de interiores, a fase de legalização junto aos órgãos responsáveis ou levantamento cadastral.
Sendo esta uma questão considerada fácil — fato que não se pode negar —, vale refletir sobre a ausência de uma terminologia adequada e específica para a área de Interiores, no que diz respeito à metodologia de projeto.
Valendo -se, como tem feito, de técnicas e terminologia de áreas afins, e sendo o Design de Interiores praticado também por profissionais de área correlatas ( que de \textbf{certo} modo contribuem para essa apropriação ), pode -se imaginar a imperiosa necessidade de estabelecer um \textbf{certo} padrão específico para a área, o qual fosse mais bem assimilado pelos seus egressos. 6.3.3 Considerações Finais As questões discursivas propostas têm total pertinência à área de Design de Interiores, e seus conteúdos correspondem a matéria que se supõe ser domínio de qualquer egresso com esta formação.
Os resultados menos satisfatórios parecem \textbf{estar}, surpreendentemente, com a Questão 5 - referente justamente à metodologia e prática de projeto, assunto de relevância para o campo de Interiores.
Nesse quesito, é importante observar que a ausência de uso de uma linguagem padronizada na área, que se serve de terminologia de disciplinas afins, dificulta a compreensão e gera equívocos.
A Questão mais fácil, a de número 3, de \textbf{ordem} mais técnica, obteve resultados um pouco melhores.
De modo geral, a dificuldade de expressão entre os estudantes permanece como um dos agravantes aos resultados insuficientes.
A própria natureza da prova induz a uma \textbf{certa} displicência observada na elaboração das respostas.
Por vezes há indícios de um relativo domínio sobre o assunto sem que se obtenha, do estudante, o cuidado necessário para o atendimento de todos os quesitos solicitados nos enunciados.
É oportuno salientar que a pontuação obtida pode \textbf{estar} associada ainda à uma \textbf{certa} ausência de planejamento na organização das respostas, o que ocasiona lapsos e esquecimentos e consequente falta de pontuação.
Mas não se pode deixar de afirmar o número demasiado alto de erros ortográficos e de concordância nominal e verbal, equívocos praticados pelos estudantes de todas as áreas.
O que pode ser um indicativo relevante para o campo de Design de Interiores é a predominância de uso inadequado da terminologia específica pertinente, tanto nas questões mais técnicas quanto nas de cunho cultural, o que evidencia que, talvez mais até do que o conteúdo, haja grande dificuldade, pelos estudantes, de incorporação dessas expressões ao seu vocabulário.

% matched lemmas: cada, certo, detalhamento, esquema, estar, ordem, torno
\end{textsample}
